\makeatletter
% Add the ICLR 2026 template directory to LaTeX's search path.
% This lets \usepackage{iclr2026_conference} work without copying files.
\def\input@path{{../../doc/iclr-template/master/iclr2026/}}
\makeatother

\documentclass{article}

% ICLR template (located in ../doc/templates/iclr2026/)
\usepackage{iclr2026_conference}
\usepackage{times}

% Math & symbols
\usepackage{amsmath,amsfonts,amssymb,amsthm}

% Figures & tables
\usepackage{graphicx}
\usepackage{booktabs}
\usepackage{subcaption}

% Formatting
\usepackage{microtype}
\usepackage{hyperref}
\usepackage{xcolor}

% Reduce underfull vbox warnings by allowing variable page height.
\raggedbottom

% Theorem-like environments (used in \input files).
\newtheorem{theorem}{Theorem}
\newtheorem{lemma}{Lemma}
\newtheorem{proposition}{Proposition}
\newtheorem{corollary}{Corollary}
\newtheorem{assumption}{Assumption}

\title{The Curse and Blessing of Bias in Large Language Model Training}

\author{
  Anonymous Authors\\
  Anonymous Institutions\\
  \texttt{anonymous@iclr.cc}
}

% Uncomment for camera-ready
% \iclrfinalcopy

\begin{document}

\maketitle

\begin{abstract}

Training of large language models (LLMs) is governed by the interaction between
forward activations and backward gradient outputs at every layer.
While prior work has emphasized spectral anisotropy in these tensors,
we show that a more elementary structural component --- 
\emph{activation bias}, a coherent mean embedded in forward representations ---
plays a central role in shaping training geometry.

Across layers and training stages, activation means exhibit
noise-normalized magnitudes exceeding $10\times$--$40\times$
the expected stochastic scale, whereas gradient-output means remain at noise level.
Furthermore, among the top $0.1\%$ largest activation entries,
$85\%$--$90\%$ of squared magnitude is attributable to the mean component,
with spike variance contributing modestly and long-tail components contributing
negligibly.
These results establish \emph{mean-dominated outliers} as a structural property
of LLM activations rather than an incidental consequence of heavy-tailed variance.

We formalize this phenomenon and show that bias propagates multiplicatively
through the outer-product gradient structure, influencing both the spectral
geometry of parameter updates and the elementwise dynamic range encountered
during training.
A simple intervention --- activation mean removal ---
isolates and eliminates the dominant coherent component while preserving
higher-order semantic covariance.
Additional spike whitening provides incremental conditioning,
but is not the primary source of improvement.

To illustrate the broader implications of this structural insight,
we evaluate bias mitigation in two downstream settings:
FP4 low-bit training and adaptive optimization.
In both cases, removing activation bias yields consistent improvements,
explaining a substantial portion of the gains previously attributed to
computationally intensive spectral methods.

These findings reposition bias as a central geometric feature of
large language model training.
While bias contributes to dynamic-range inflation and conditioning pathologies,
its structured coherence also enables efficient mitigation,
revealing both the curse and the blessing of bias.

\end{abstract}

\section{Introduction}

Large language models are trained on massive human language corpora,
which induce a characteristic geometry in learned representations.
Across layers and training stages, activations, gradients, and model parameters
exhibit pronounced anisotropy:
a small number of directions concentrate disproportionate energy
(strong spectral spikes),
while the remaining directions form a broad long tail
\cite{ethayarajh2019contextual,mu2018all,timkey2021rogue,kovaleva2021bertbusters,machina2024anisotropy}.

This anisotropy is not accidental.
It reflects the hierarchical structure of human language itself.
At every level of representation—tokens, phrases, clauses, and full sentences—
there exist shared semantic components that are common across tokens within a context.
When contextual aggregation mechanisms such as attention and nonlinear transformations
extract these shared semantics, they naturally induce coherent mean directions
in the representation space.
Consequently, mean bias is not confined to the embedding layer;
it re-emerges across layers as models progressively capture higher-level
semantic regularities.

While such structure is linguistically natural,
it becomes a computational liability during large-scale training.
Anisotropy inflates elementwise dynamic range,
distorts adaptive second-moment estimates
\cite{kingma2014adam,loshchilov2017adamw,spectra2026},
and compresses long-tail variation in low-bit regimes
\cite{xiao2023smoothquant,frantar2022gptq,lin2024awq,li2024svdquant,ashkboos2024quarot}.
In high-dimensional spaces, even moderate coherent bias produces
a consistent directional drift across thousands of coordinates,
amplifying outliers and destabilizing numerical precision.
Thus anisotropy constitutes a structural \emph{curse}
for efficient and stable large-scale training.

In this work, we perform a systematic layer-wise investigation
of anisotropy across activations, gradient outputs, and model parameters.
We uncover a structural asymmetry:
\emph{mean bias is a dominant component of activation anisotropy across layers,
whereas gradient outputs exhibit no comparable coherent mean}.
Moreover, although model parameters display spectral spikes,
they do not accumulate persistent mean offsets.
This establishes activation mean bias as the principal forward-path driver
of anisotropic geometry.

Importantly, we show that anisotropy can be decomposed into two components:
(i) a coherent mean bias, and
(ii) a residual low-rank spike.
The mean bias is the primary source of extreme elementwise outliers
and adaptive second-moment inflation,
while the residual spike plays a secondary role.
Thus the dominant instability mechanism is not generic high-dimensional variance,
but a structured rank-one component that propagates through residual connections
across depth.

Tracing this structure backward,
we identify its statistical origin in the skewed token distribution of natural language.
Embedding weights develop frequency-aligned structure
\cite{mikolov2013efficient,mikolov2013distributed,arora2017simple},
and subsequent contextual aggregation amplifies shared semantics at progressively
higher abstraction levels.
Mean bias therefore reflects the hierarchical organization of world knowledge
encoded in language rather than an architectural accident.

Crucially, this structure is not merely harmful—it is also exploitable.
Because the dominant instability arises from a rank-one component,
it can be isolated and removed with a single mean-subtraction operation.
The remaining low-rank spike can, if necessary,
be whitened once at the embedding activation level.
Both operations are hardware-friendly and substantially cheaper
than repeated layer-wise SVD corrections.
In this sense, the same structure that creates instability
also enables efficient mitigation—the \emph{blessing} of bias.

We validate this bias-centric interpretation in two downstream settings.
In low-bit quantization, mean removal significantly reduces dynamic-range inflation
\cite{xiao2023smoothquant,wei2023outlier,ashkboos2024quarot}.
In adaptive optimization, it contracts extreme second-moment tails,
improving effective learning-rate allocation
\cite{kingma2014adam,spectra2026}.
These results suggest that addressing bias at its source
can replace or complement spectral interventions
with minimal computational overhead.

\paragraph{Contributions.}
\begin{enumerate}[leftmargin=*,itemsep=2pt]
    \item We show that activation mean bias emerges across layers as a structural consequence of hierarchical semantic aggregation.
    \item We demonstrate that mean-dominated components account for most extreme activation magnitudes and second-moment inflation.
    \item We establish a forward-path asymmetry: gradient outputs remain mean-neutral while activations accumulate coherent bias.
    \item We propose source-level mean removal and optional spike whitening as efficient spectral conditioning mechanisms.
    \item We demonstrate improvements in low-bit quantization and adaptive optimization.
\end{enumerate}
\section{Related Work}
\label{sec:related}

\subsection{Frequency Bias and Subsampling in Word2Vec (2013)}

The earliest large-scale embedding systems already encountered statistical
dominance of frequent tokens.
In the original Word2Vec formulation~\cite{mikolov2013efficient},
high-frequency words were observed to contribute disproportionately
to training updates.
To mitigate this, subsampling was introduced:
frequent words are randomly discarded with probability increasing
in their corpus frequency.

The stated motivation was computational efficiency.
However, from a geometric perspective,
subsampling implicitly suppresses a dominant first-order component
associated with corpus frequency.
This can be interpreted as an early heuristic attempt to control
frequency-driven anisotropy in the embedding space.

Importantly, Word2Vec did not explicitly analyze the embedding mean
or its contribution to spectral structure.
Bias mitigation was treated as an optimization trick,
rather than as a structural property of the learned representation.

\subsection{Common Direction Removal in Sentence Embeddings (SIF, 2017)}

Arora et al.~\cite{arora2017simple} proposed Smooth Inverse Frequency (SIF)
sentence embeddings, observing that averaged word vectors tend to
be dominated by a common direction.
They remove the top principal component after sentence averaging,
demonstrating improved downstream performance.

Their interpretation is that frequent words induce a dominant
latent direction that does not encode semantic content.
This work introduced an explicit spectral removal step
(post-hoc PCA subtraction) and formalized a generative model
in which a common discourse vector explains the phenomenon.

While SIF identified and removed a dominant direction,
it did not distinguish between mean-induced structure
and higher-order covariance anisotropy,
nor did it analyze the impact on training dynamics or optimization.

\subsection{Anisotropy in Word Embeddings (All-but-the-Top, 2018)}

Mu and Viswanath~\cite{mu2018all} systematically analyzed anisotropy
in pretrained word embeddings.
They observed that embeddings concentrate in a narrow cone
and that removing the mean and the top principal components
significantly improves isotropy and performance.

The All-but-the-Top (ABTT) procedure subtracts the mean
and removes several leading singular vectors.
This explicitly connects anisotropy to low-rank structure
and demonstrates that dominant spectral directions
can be detrimental to representation quality.

However, ABTT operates as a post-processing technique
on static embeddings.
It does not analyze the causal origin of the dominant component,
nor does it connect mean-dominated structure to
optimization instability or quantization pathologies.

\subsection{Positioning of the Present Work}

The prior literature establishes three important facts:

\begin{itemize}
\item High-frequency tokens induce dominant directions.
\item Embedding spaces are strongly anisotropic.
\item Removing leading spectral components improves representation quality.
\end{itemize}

Our work differs in scope and focus.

First, we identify that in modern transformer training,
the dominant first-order structure is specifically an
\emph{activation mean bias} emerging from embedding statistics
and token frequency imbalance.

Second, we show that this bias is not merely a representation artifact,
but propagates multiplicatively through gradient outer products,
inflating optimizer second moments and quantization scales.

Third, we provide a unifying causal analysis linking
frequency imbalance $\rightarrow$ embedding bias
$\rightarrow$ activation anisotropy
$\rightarrow$ gradient-scale distortion.

Thus, while earlier work removed dominant directions
to improve embedding quality,
we show that bias is a central structural driver
of multiple downstream training phenomena
in large language models.:w

\section{Empirical Discovery: Mean-Dominated Geometry}
\label{sec:empirical}

This section presents the central empirical finding of the paper:
activation tensors in large language models contain a dense, coherent mean
component that dominates extreme element-wise outliers and materially
contributes to spectral anisotropy.
We contrast this behavior with gradient-output tensors and establish that
bias is activation-specific.

\subsection{First-Order Structure of Training Signals}
\label{subsec:first_order_asymmetry}

Training of large language models is driven by two primary data objects at each layer:
the forward activations $x_i \in \mathbb{R}^H$ and the backward gradient outputs
$\delta_i \in \mathbb{R}^H$.
These two tensors jointly determine the parameter update through outer-product
structure
\[
g_i = \delta_i x_i^\top.
\]
Any systematic structure present in either $x_i$ or $\delta_i$
directly propagates into the gradient matrix and therefore into the model parameters.

We begin by examining their first-order statistics.
For a batch of $N$ tokens, define the empirical mean
\[
\mu_z = \frac{1}{N} \sum_{i=1}^N z_i,
\qquad z_i \in \{x_i, \delta_i\}.
\]

To determine whether $\mu_z$ reflects coherent structure or merely finite-sample
fluctuation, we compare its magnitude to the expected size of a noise-level mean.
If $\{z_i\}$ were independent zero-mean samples with variance
$\mathbb{E}\|z_i\|_2^2$, then by standard concentration,
\[
\|\mu_z\|_2 \approx \sqrt{\frac{\mathbb{E}\|z_i\|_2^2}{N}}.
\]
This motivates the noise-normalized mean magnitude
\[
R_z
=
\frac{\|\mu_z\|_2}
{\sqrt{\mathbb{E}\|z_i\|_2^2 / N}}.
\]

Interpretation:
\begin{itemize}
\item $R_z \approx 1$ indicates that the empirical mean is consistent with
finite-sample noise.
\item $R_z \gg 1$ indicates coherent, batch-aligned structure.
\end{itemize}

\paragraph{Empirical measurement.}

Across layers (3 and 21) and training steps (10k and 180k),
we observe a pronounced asymmetry:

\begin{itemize}
\item Activations exhibit large coherent means,
      with $R_x$ ranging from approximately $17$ to $41$.
\item Gradient outputs remain near noise level,
      with $R_\delta$ ranging from approximately $1.05$ to $1.19$.
\end{itemize}

Figure~\ref{fig:activation_grad_R} reports these values.

\begin{figure}[t]
\centering
\includegraphics[width=0.7\linewidth]{figures/sec2_1_R_bar.png}
\caption{
Noise-normalized mean magnitude $R_z$ for activations and gradient outputs.
Activations exhibit strong coherent first-order structure ($R_x \gg 1$),
whereas gradient outputs remain consistent with noise ($R_\delta \approx 1$).
}
\label{fig:activation_grad_R}
\end{figure}

\paragraph{Observation.}

The data reveal a clear first-order asymmetry:
coherent mean structure is present in activations but absent in gradient outputs.
Since gradients are constructed as outer products $g_i = \delta_i x_i^\top$,
this implies that coherent first-order bias in training must originate
from the activation side.

This asymmetry motivates the subsequent analysis.
If coherent bias dominates the forward signal but not the backward signal,
then anisotropy observed in gradients and optimizer statistics
must be driven primarily by activation structure.

\subsection{Outliers Are Dominated by the Activation Mean}
\label{subsec:outlier_mean_dominance}

We now identify the origin of extreme elementwise magnitudes in training tensors.
While spectral analyses describe variance structure at the matrix level,
low-bit quantization and several optimizer pathologies are triggered by
\emph{entrywise} extremes.
We therefore ask:

\begin{quote}
\emph{Among the largest-magnitude entries, which component contributes most:
mean bias, spike variance, or long-tail variance?}
\end{quote}

\paragraph{Decomposition.}
For each matrix $X \in \mathbb{R}^{N\times H}$ (activation or gradient output),
we decompose
\[
X = M + X_{\text{spike}} + X_{\text{tail}},
\]
where $M=\mathbf{1}\mu^\top$ is the broadcasted mean component,
$X_{\text{spike}}$ is the projection of the centered matrix onto the top-$k$
right singular vectors (we use $k=1\%$ of $H$),
and $X_{\text{tail}}$ is the remaining component.

\paragraph{Outlier set and contribution shares.}
We define outliers as the top $0.1\%$ entries ranked by raw magnitude $|X_{ij}|$.
For each outlier entry, we compute the squared contribution shares
\[
\rho_{\text{mean}} = \frac{M_{ij}^2}{X_{ij}^2},\qquad
\rho_{\text{spike}} = \frac{(X_{\text{spike}})_{ij}^2}{X_{ij}^2},\qquad
\rho_{\text{tail}} = \frac{(X_{\text{tail}})_{ij}^2}{X_{ij}^2}.
\]
The resulting distributions of $\rho_{\text{mean}},\rho_{\text{spike}},\rho_{\text{tail}}$
provide a direct elementwise attribution of extreme magnitudes.

\paragraph{Results across layers and training steps.}
Figure~\ref{fig:outlier_contrib_grid} shows the outlier contribution histograms
for activations (left column) and gradient outputs (right column),
across two layers (3 and 21) and two training steps (10k and 180k).

Two consistent patterns emerge:
\begin{itemize}
    \item \textbf{Activations: mean-dominated outliers.}
    Activation outliers concentrate strongly near $\rho_{\text{mean}}\approx 1$,
    indicating that the mean component accounts for the overwhelming majority of
    squared magnitude in the most extreme entries.
    Spike contributions are secondary and tail contributions are negligible.
    \item \textbf{Gradient outputs: variance-dominated outliers.}
    Gradient-output outliers show negligible mean contribution
    ($\rho_{\text{mean}}\approx 0$ for nearly all outliers) and are instead
    dominated by variance components (spike and/or tail).
\end{itemize}

\paragraph{Implication.}
This establishes the central empirical fact behind this paper:
\begin{quote}
\emph{Elementwise outliers in activations are overwhelmingly bias-driven, whereas
gradient-output outliers are variance-driven.}
\end{quote}
Consequently, mean subtraction targets the primary source of quantization-relevant
activation outliers, while downstream spectral conditioning addresses residual
variance anisotropy.

\begin{figure*}[t]
\centering

% ================= Row 1 =================
\begin{subfigure}[t]{0.23\linewidth}
  \centering
  \includegraphics[width=\linewidth]{figures/layer3_step010000_activation.pt_outlier_hist.png}
  \caption{L3 10k (Act)}
\end{subfigure}
\hfill
\begin{subfigure}[t]{0.23\linewidth}
  \centering
  \includegraphics[width=\linewidth]{figures/layer3_step010000_grad_output.pt_outlier_hist.png}
  \caption{L3 10k (Grad)}
\end{subfigure}
\hfill
\begin{subfigure}[t]{0.23\linewidth}
  \centering
  \includegraphics[width=\linewidth]{figures/layer3_step180000_activation.pt_outlier_hist.png}
  \caption{L3 180k (Act)}
\end{subfigure}
\hfill
\begin{subfigure}[t]{0.23\linewidth}
  \centering
  \includegraphics[width=\linewidth]{figures/layer3_step180000_grad_output.pt_outlier_hist.png}
  \caption{L3 180k (Grad)}
\end{subfigure}

\vspace{0.6em}

% ================= Row 2 =================
\begin{subfigure}[t]{0.23\linewidth}
  \centering
  \includegraphics[width=\linewidth]{figures/layer21_step010000_activation.pt_outlier_hist.png}
  \caption{L21 10k (Act)}
\end{subfigure}
\hfill
\begin{subfigure}[t]{0.23\linewidth}
  \centering
  \includegraphics[width=\linewidth]{figures/layer21_step010000_grad_output.pt_outlier_hist.png}
  \caption{L21 10k (Grad)}
\end{subfigure}
\hfill
\begin{subfigure}[t]{0.23\linewidth}
  \centering
  \includegraphics[width=\linewidth]{figures/layer21_step180000_activation.pt_outlier_hist.png}
  \caption{L21 180k (Act)}
\end{subfigure}
\hfill
\begin{subfigure}[t]{0.23\linewidth}
  \centering
  \includegraphics[width=\linewidth]{figures/layer21_step180000_grad_output.pt_outlier_hist.png}
  \caption{L21 180k (Grad)}
\end{subfigure}

\caption{
Outlier contribution decomposition across layers and training steps.
Each panel shows squared contribution shares for the top $0.1\%$ largest entries.
Across all settings, activation outliers are overwhelmingly mean-dominated,
whereas gradient-output outliers are variance-dominated.
}
\label{fig:outlier_contrib_grid}

\end{figure*}


\subsection{Spectral Role of Bias}
\label{subsec:spectral_role_bias}

Having established that activation bias dominates extreme entries, we now
investigate how this bias manifests in the singular value spectrum.

Let $X \in \mathbb{R}^{N \times H}$ denote the activation matrix.
We compute the singular value spectrum of both:

\begin{itemize}
\item the raw activation matrix $X$,
\item the mean-centered activation matrix $X_c = X - \mathbf{1}\mu^\top$.
\end{itemize}

Our goal is to determine:
\begin{enumerate}
\item how much of the spectral spike is attributable to bias,
\item whether bias accounts for the entire spike,
\item how this behavior varies across layers and training stages.
\end{enumerate}

\paragraph{Experimental setup.}
For each of four configurations
(layer 3 and layer 21 at 10k and 180k steps),
we compute:

\begin{itemize}
\item singular value spectrum (log-log scale),
\item cumulative spectral energy.
\end{itemize}

The spike rank $k$ is defined as $1\%$ of the feature dimension.

\paragraph{Results.}

\begin{figure}[t]
\centering
\begin{subfigure}[t]{0.24\linewidth}
\includegraphics[width=\linewidth]{figures/layer3_10k_spectrum_raw_vs_centered.png}
\caption{Layer 3, 10k (spectrum)}
\end{subfigure}
\begin{subfigure}[t]{0.24\linewidth}
\includegraphics[width=\linewidth]{figures/layer3_10k_cum_energy_raw_vs_centered.png}
\caption{Layer 3, 10k (energy)}
\end{subfigure}
\begin{subfigure}[t]{0.24\linewidth}
\includegraphics[width=\linewidth]{figures/layer3_180k_spectrum_raw_vs_centered.png}
\caption{Layer 3, 180k (spectrum)}
\end{subfigure}
\begin{subfigure}[t]{0.24\linewidth}
\includegraphics[width=\linewidth]{figures/layer3_180k_cum_energy_raw_vs_centered.png}
\caption{Layer 3, 180k (energy)}
\end{subfigure}

\vspace{0.5em}

\begin{subfigure}[t]{0.24\linewidth}
\includegraphics[width=\linewidth]{figures/layer21_10k_spectrum_raw_vs_centered.png}
\caption{Layer 21, 10k (spectrum)}
\end{subfigure}
\begin{subfigure}[t]{0.24\linewidth}
\includegraphics[width=\linewidth]{figures/layer21_10k_cum_energy_raw_vs_centered.png}
\caption{Layer 21, 10k (energy)}
\end{subfigure}
\begin{subfigure}[t]{0.24\linewidth}
\includegraphics[width=\linewidth]{figures/layer21_180k_spectrum_raw_vs_centered.png}
\caption{Layer 21, 180k (spectrum)}
\end{subfigure}
\begin{subfigure}[t]{0.24\linewidth}
\includegraphics[width=\linewidth]{figures/layer21_180k_cum_energy_raw_vs_centered.png}
\caption{Layer 21, 180k (energy)}
\end{subfigure}

\caption{
Singular spectra and cumulative energy before and after mean removal.
Mean removal significantly suppresses the leading singular component
in early layers, while leaving the long tail largely unchanged.
}
\label{fig:spectral_bias_role}
\end{figure}

\paragraph{Observation 1: Bias strongly contributes to spike energy in early layers.}

In layer 3 at 10k steps, mean removal dramatically reduces the leading singular
value and shifts the cumulative energy curve downward for small $k$.
This confirms that a substantial portion of the spectral spike is due to
the coherent mean component.

\paragraph{Observation 2: Bias is not the entire spike.}

Even after mean removal, a clear low-rank structure remains.
The spike does not vanish entirely.
This indicates that anisotropy contains two components:

\begin{itemize}
\item a dense coherent mean component,
\item structured low-rank variance.
\end{itemize}

\paragraph{Observation 3: Bias influence weakens in deeper layers.}

For layer 21 at 10k, raw and centered spectra are nearly identical.
This indicates that deep-layer representations contain substantially less
mean-driven anisotropy.

\paragraph{Observation 4: Bias evolves during training.}

Comparing 10k and 180k steps shows that bias-driven spike strength
changes during training.
Early stages exhibit stronger mean-dominated structure.

\paragraph{Conclusion.}

The spectral spike decomposes into:

\begin{equation}
\text{Spike} = \text{Dense coherent mean} + \text{Low-rank structured variance}.
\end{equation}

Mean removal eliminates the first term but leaves the second intact.
Therefore:

\begin{itemize}
\item Bias is a structured component inside the spike,
\item but it is not synonymous with the entire spike,
\item and spike variance persists after mean removal.
\end{itemize}

This refined understanding will be essential in Section~\ref{sec:grad_propagation},
where we analyze how bias propagates into gradient updates.

\subsection{Bias Propagation into Gradient Spectral Geometry}
\label{subsec:grad_spectrum_bias}

Having established that activation tensors contain a statistically significant
dense mean component (Section~2.1) and that this mean dominates extreme
activation outliers (Section~2.2), we now examine how this bias propagates into
the gradient matrix that directly drives parameter updates.

\paragraph{Gradient structure.}
For a linear layer, the per-token gradient has outer-product structure
\[
g_i = \delta_i x_i^\top,
\qquad
G = \frac{1}{N}\sum_i g_i = \frac{1}{N}\Delta^\top X.
\]
Although gradient outputs $\delta_i$ exhibit no significant mean bias
(Section~2.1), the activation matrix $X$ does.
Therefore, even if $\Delta$ is approximately zero-mean, the product
$\Delta^\top X$ may inherit structured anisotropy through $X$.

To isolate this effect, we compare:
\[
G_{\text{raw}} = \frac{1}{N}\Delta^\top X,
\qquad
G_{\text{centered}} = \frac{1}{N}\Delta^\top (X - \mu_x).
\]

\paragraph{Experimental protocol.}
For each layer and training step, we compute:
\begin{itemize}
\item the singular value spectrum of $G_{\text{raw}}$,
\item the singular value spectrum of $G_{\text{centered}}$,
\item the cumulative spectral energy distribution.
\end{itemize}
All spectra are plotted on log–log scale to reveal spike–tail separation.

\begin{figure}[t]
\centering

% Row 1
\begin{subfigure}[t]{0.24\linewidth}
\includegraphics[width=\linewidth]{figures/layer3_10k_G_spectrum_raw_vs_centered.png}
\caption{Layer 3, 10k (spectrum)}
\end{subfigure}
\begin{subfigure}[t]{0.24\linewidth}
\includegraphics[width=\linewidth]{figures/layer3_10k_G_cum_energy_raw_vs_centered.png}
\caption{Layer 3, 10k (cumulative)}
\end{subfigure}
\begin{subfigure}[t]{0.24\linewidth}
\includegraphics[width=\linewidth]{figures/layer3_180k_G_spectrum_raw_vs_centered.png}
\caption{Layer 3, 180k (spectrum)}
\end{subfigure}
\begin{subfigure}[t]{0.24\linewidth}
\includegraphics[width=\linewidth]{figures/layer3_180k_G_cum_energy_raw_vs_centered.png}
\caption{Layer 3, 180k (cumulative)}
\end{subfigure}

\vspace{0.5em}

% Row 2
\begin{subfigure}[t]{0.24\linewidth}
\includegraphics[width=\linewidth]{figures/layer21_10k_G_spectrum_raw_vs_centered.png}
\caption{Layer 21, 10k (spectrum)}
\end{subfigure}
\begin{subfigure}[t]{0.24\linewidth}
\includegraphics[width=\linewidth]{figures/layer21_10k_G_cum_energy_raw_vs_centered.png}
\caption{Layer 21, 10k (cumulative)}
\end{subfigure}
\begin{subfigure}[t]{0.24\linewidth}
\includegraphics[width=\linewidth]{figures/layer21_180k_G_spectrum_raw_vs_centered.png}
\caption{Layer 21, 180k (spectrum)}
\end{subfigure}
\begin{subfigure}[t]{0.24\linewidth}
\includegraphics[width=\linewidth]{figures/layer21_180k_G_cum_energy_raw_vs_centered.png}
\caption{Layer 21, 180k (cumulative)}
\end{subfigure}

\caption{Gradient singular spectra before and after activation mean removal.
Across layers and training stages, removing the activation mean consistently
reduces the dominance of the leading singular directions and shifts spectral
energy toward the long tail.}
\label{fig:grad_spectrum_bias}
\end{figure}

\paragraph{Observations.}

Across all layers and training stages, three consistent phenomena emerge:

\begin{enumerate}
\item \textbf{Raw gradients exhibit spike–tail structure.}
Despite the absence of mean bias in gradient outputs $\delta_i$,
the gradient matrix $G$ displays strong anisotropy, confirming that
activation geometry propagates into the optimizer.

\item \textbf{Mean removal weakens spike dominance.}
After centering activations, the leading singular values shrink,
and cumulative energy curves shift rightward.
The spike becomes less dominant relative to the tail.

\item \textbf{The spike does not vanish.}
Even after mean removal, a structured low-rank component remains.
Thus, the spectral spike consists of two parts:
\[
\text{Spike} = \text{Dense mean component} + \text{Structured covariance component}.
\]
Mean removal eliminates the dense bias term but leaves the structured
covariance component intact.
\end{enumerate}

\paragraph{Interpretation.}

These results establish a key structural asymmetry:

\begin{itemize}
\item Activation bias directly inflates gradient anisotropy.
\item Removing activation bias reduces gradient spike energy
without affecting gradient-output statistics.
\item The remaining spike reflects genuine covariance structure rather
than trivial mean effects.
\end{itemize}

This separation is crucial for understanding optimization behavior:
activation bias acts as a low-rank energy amplifier that artificially
inflates second-moment statistics, while structured covariance encodes
learning-relevant variability.

In Section~2.4B we quantify how this bias inflation impacts adaptive
second-moment optimizers.

\subsection{Parameter Matrices Do Not Accumulate Bias}
\label{subsec:parameter_bias}

Previous sections
established that activation tensors exhibit strong coherent mean bias,
and that this bias propagates multiplicatively into gradient matrices.
We now examine whether such bias accumulates in the model parameters
themselves during training.

\paragraph{Measurement protocol.}
For representative feed-forward layers (Layer~3 and Layer~21),
we analyze weight matrices
$W \in \mathbb{R}^{3072 \times 1024}$
at early (10k) and later (180k) training steps.
We compute both column-wise and row-wise mean components and evaluate:

\begin{enumerate}
\item Noise-normalized mean magnitude
\[
R = \frac{\|\mu\|_2}{\sqrt{\mathbb{E}\|W\|_F^2 / d}},
\]
where $d$ denotes the corresponding dimensionality.
\item Mean energy fraction
\[
\rho = \frac{\|\mathbf{1}\mu^\top\|_F^2}{\|W\|_F^2}.
\]
\end{enumerate}

\begin{table}[t]
\centering
\caption{Bias statistics of model parameters across layers and training stages.}
\label{tab:weight_bias}
\begin{tabular}{lcccc}
\toprule
Layer & Step & $R_{\text{col}}$ & $\rho_{\text{col}}$ & $\rho_{\text{row}}$ \\
\midrule
Layer 3  & 10k   & 0.970 & $3.06\times10^{-4}$ & $9.85\times10^{-4}$ \\
Layer 3  & 180k  & 0.988 & $3.18\times10^{-4}$ & $9.89\times10^{-4}$ \\
Layer 21 & 10k   & 0.999 & $3.25\times10^{-4}$ & $9.81\times10^{-4}$ \\
Layer 21 & 180k  & 1.002 & $3.27\times10^{-4}$ & $9.65\times10^{-4}$ \\
\bottomrule
\end{tabular}
\end{table}

\paragraph{Observations.}

Across all layers and training stages:

\begin{itemize}
\item Noise-normalized mean ratios remain at $R \approx 1$,
indicating that parameter means are indistinguishable from
finite-sample noise.
\item Mean energy fractions remain below $0.1\%$ of total Frobenius energy.
\item No systematic growth of parameter bias is observed over training.
\end{itemize}

In contrast, activation tensors exhibit
$R_x \in [17, 41]$ (Section~\ref{subsec:mean_asymmetry}),
demonstrating that coherent bias is specific to intermediate
representations rather than learned parameters.

\paragraph{Implication.}
These results establish that mean-dominated geometry is not
a persistent property of model weights.
Bias originates in activations, influences gradients transiently,
but does not accumulate into parameter matrices.
This distinction localizes the geometric source of bias-driven
anisotropy to the forward representation space.

\subsection{Bias Originates at the Embedding Layer}
\label{subsec:embedding_bias}

To determine whether activation bias is induced by deeper layers or
already present at the input stage, we analyze the embedding-layer
activations, i.e., the token representations fed into the transformer
stack.

Let $x_{b,t} \in \mathbb{R}^H$ denote the embedding of token $t$ in
sequence $b$, and let $X \in \mathbb{R}^{(B L)\times H}$ be the
flattened batch matrix.
We measure bias magnitude using the same noise-normalized statistic
introduced in Section~2.1:
\[
R \;=\;
\frac{\|\mu\|_2}
{\sqrt{\mathbb{E}\|x\|_2^2 / N}},
\qquad
\mu = \frac{1}{N}\sum_{i=1}^N x_i,
\]
where $N = B L$.

\paragraph{Batch-level bias.}
For the embedding layer we observe:
\[
R_{\text{batch}} \approx 24.0 \text{ (step 10k)},
\qquad
R_{\text{batch}} \approx 25.6 \text{ (step 180k)}.
\]
Since $R\approx 1$ corresponds to noise-level mean,
these values indicate that embedding activations possess a
highly significant coherent mean component.

\paragraph{Per-sequence bias.}
We further compute sequence-level means
\[
\mu_b = \frac{1}{L} \sum_{t=1}^L x_{b,t}
\]
for each of the $B=8$ sequences in the batch.

The average per-sequence noise-normalized ratio is:
\[
\mathbb{E}_b[R_{\text{seq}}] \approx 8.8 \text{ (10k)},
\qquad
\mathbb{E}_b[R_{\text{seq}}] \approx 9.3 \text{ (180k)}.
\]

Thus, bias is already strong within individual sequences,
and is not merely a batching artifact.

\paragraph{Cross-sequence coherence.}
To quantify directional consistency, we measure the cosine
similarity between each sequence mean and the global batch mean:
\[
\cos(\mu_b, \mu_{\text{batch}}).
\]

We obtain:
\[
\text{mean cosine} \approx 0.968 \text{ (10k)},
\qquad
\text{mean cosine} \approx 0.976 \text{ (180k)}.
\]

These near-unity values indicate that sequence-level biases
are almost perfectly aligned across sequences.

\paragraph{Implications.}
These results establish three structural facts:

\begin{enumerate}
\item The embedding layer already contains a strong coherent bias.
\item This bias is present within individual sequences.
\item The bias direction is globally aligned across sequences.
\end{enumerate}

Therefore, activation bias is not created by deeper nonlinear
transformations; rather, it originates at the embedding stage
and propagates through the network.

This observation is critical:
the geometric anisotropy observed in deeper layers is seeded
by a coherent input bias rather than emerging purely from
layerwise covariance structure.

\subsubsection{Causal Origin of Activation Bias: Token Frequency Structure}
\label{subsec:token_frequency_bias}

The preceding subsections established that activation bias is
(i) statistically significant,
(ii) responsible for extreme outliers,
and (iii) absent in gradient-output tensors and learned weight matrices.
We now identify its structural origin.

\paragraph{Frequency-weighted embedding mean.}

Let $E \in \mathbb{R}^{V \times H}$ denote the learned embedding table,
and let $p(v)$ denote the empirical token frequency in a batch.
The activation mean induced by token embeddings is

\[
\mu_{\text{freq}}
=
\sum_{v=1}^{V} p(v)\, E[v].
\]

This expression shows that activation bias is not arbitrary:
it is the frequency-weighted average of embedding vectors.
If high-frequency tokens share coherent geometric structure,
$\mu_{\text{freq}}$ becomes large.

\paragraph{Concentration of bias contribution.}

Using $1{,}024 \times 1{,}024$ sampled tokens
($1{,}048{,}576$ total tokens),
we measured how much of $\|\mu_{\text{freq}}\|_2$
is explained by the top-$K$ most frequent tokens.

Results show:

\begin{itemize}
\item Top 10 tokens explain $\approx 62\%$ of $\|\mu_{\text{freq}}\|$.
\item Top 20 tokens explain $\approx 83\%$.
\item Top 50 tokens explain $\approx 110\%$ (due to constructive alignment).
\item Top 200 tokens explain $\approx 132\%$.
\end{itemize}

This demonstrates that activation bias is highly concentrated
in a small subset of frequent tokens.

\paragraph{Identity of dominant tokens.}

Table~\ref{tab:top_tokens_bias} lists the top contributing tokens
by absolute bias contribution.
These tokens are overwhelmingly function words and structural markers.

\begin{table}[t]
\centering
\caption{Top tokens contributing to embedding-induced activation bias (Layer 3, 180k step).}
\label{tab:top_tokens_bias}
\begin{tabular}{rll}
\toprule
Rank & Token ID & Decoded Token \\
\midrule
1  & 151643 & \texttt{<|endoftext|>} \\
2  & 279    & \texttt{the} \\
3  & 11     & \texttt{,} \\
4  & 13     & \texttt{to} \\
5  & 311    & \texttt{and} \\
6  & 323    & \texttt{of} \\
7  & 315    & \texttt{a} \\
8  & 264    & \texttt{in} \\
9  & 220    & \texttt{is} \\
10 & ...    & \texttt{that} \\
\midrule
\multicolumn{3}{c}{Additional high contributors: \texttt{you, for, it, I, on, are, with, be, as, 's, have, was, or, (}} \\
\bottomrule
\end{tabular}
\end{table}

These tokens correspond to:

\begin{itemize}
\item Articles and determiners
\item Prepositions
\item Conjunctions
\item Auxiliary verbs
\item Punctuation
\item Special sequence markers
\end{itemize}

In other words, activation bias aligns strongly with grammatical structure.

\paragraph{Geometric interpretation.}

The bias vector can therefore be interpreted as a
\emph{grammar direction} in embedding space:

\[
\mu_{\text{freq}} \approx
\sum_{v \in \text{function tokens}} p(v)\, E[v].
\]

Since grammar tokens appear in nearly every sentence,
their contributions accumulate coherently across sequences,
producing a dense rank-one component in activation space.

\paragraph{Connection to earlier findings.}

This explains the empirical observations:

\begin{enumerate}
\item Activation bias magnitude ratios $R_x \gg 1$ (Section~2.1).
\item Mean-dominated outliers (Section~2.2).
\item Significant spectral energy reduction after centering (Section~2.3).
\item Gradient bias inflation via outer-product propagation (Section~2.4).
\end{enumerate}

Crucially, the bias does \emph{not} arise from
optimizer dynamics or weight accumulation;
it originates directly from corpus frequency imbalance
interacting with coherent embedding geometry.

\paragraph{Conclusion.}

Activation bias is a deterministic consequence of:

\[
\text{Zipfian token frequency}
\;\times\;
\text{coherent embedding structure}.
\]

This provides a complete causal explanation for
mean-dominated outliers in large language model training.

\subsection{Mean Removal as the First-Order Mechanism for Preserving Long-Tail Resolution}
\label{sec:mean_first_order_quantization}

Low-bit quantization allocates limited dynamic range across all activation values.
Under blockwise quantization, the quantization scale is determined by extreme
entries within each block.
Consequently, the component responsible for extreme magnitudes directly
controls representational resolution of the remaining features.

This section provides quantitative evidence that activation mean bias,
rather than residual spike variance, is the dominant driver of
long-tail degradation under low-bit quantization.

\subsubsection{Experimental Setup}

We evaluate four activation tensors:

\begin{itemize}
\item Layer 3 at step 10k (early layer, early training),
\item Layer 3 at step 180k (early layer, late training),
\item Layer 21 at step 10k (deep layer, early training),
\item Layer 21 at step 180k (deep layer, late training).
\end{itemize}

For each tensor, we:

\begin{enumerate}
\item Decompose activation into mean, spike (top 1\%), and tail components.
\item Identify the top 0.1\% largest elementwise magnitudes.
\item Compute squared contribution shares of mean, spike, and tail.
\item Perform blockwise INT4 quantization (block size = 64).
\item Measure tail signal-to-quantization-noise ratio (SQNR) at $p99.9$.
\end{enumerate}

\subsubsection{Outlier Decomposition}

Across all four settings, we observe:

\begin{itemize}
\item The mean component explains approximately $85\%$–$90\%$
of squared magnitude among the top 0.1\% largest activation entries.
\item The spike component explains roughly $5\%$–$15\%$.
\item The tail contribution is negligible.
\end{itemize}

Thus, extreme values that determine quantization scale
are overwhelmingly mean-dominated.

\subsubsection{Quantitative Impact on Tail SQNR}

We compare three conditions:

\begin{itemize}
\item Raw activation.
\item Mean removed.
\item Spike removed (mean preserved).
\end{itemize}

Representative tail SQNR results (block=64, $p99.9$):

\begin{center}
\begin{tabular}{lccc}
\toprule
Setting & Raw & Mean Removed & Spike Removed \\
\midrule
Layer 3 (10k)   & 8.7 dB  & 12.0 dB  & 10.8 dB \\
Layer 3 (180k)  & 10.6 dB & 11.9 dB  & 11.0 dB \\
Layer 21 (10k)  & 11.2 dB & 13.0 dB  & 12.0 dB \\
Layer 21 (180k) & 11.1 dB & 13.5 dB  & 12.0 dB \\
\bottomrule
\end{tabular}
\end{center}

Two consistent quantitative patterns emerge:

\begin{enumerate}
\item Removing the mean improves tail SQNR by approximately
$+2$ to $+4$ dB across settings.
\item Removing spike components while keeping the mean yields
smaller improvements.
\end{enumerate}

Therefore, the mean bias is the primary determinant
of long-tail quantization degradation.

\subsubsection{First-Order vs Second-Order Effects}

Although spike components may explain a nontrivial portion
of total spectral energy,
they contribute far less to extreme elementwise magnitudes
that determine quantization scale.

Formally:
\[
\text{Outlier Energy} \approx
\text{Mean Energy} \gg \text{Spike Energy} \gg \text{Tail Energy}.
\]

Since block quantization scale is governed by
\[
\alpha \propto \max_{i \in \text{block}} |x_i|,
\]
and these maxima are mean-dominated,
mean removal directly reduces dynamic-range inflation.

Spike whitening provides secondary refinement,
but does not address the principal source of extreme values.

\subsubsection{Conclusion}

These results establish that:

\begin{quote}
\textbf{Activation mean bias is the first-order mechanism
responsible for long-tail resolution loss under low-bit quantization.}
\end{quote}

Removing the mean alone yields the largest improvement
in tail SQNR with minimal computational cost.
Residual spike treatment is beneficial but secondary.

This finding motivates addressing bias at its source
(embedding activation)
rather than relying on repeated layer-wise spectral decompositions.

\subsection{Activation Bias and Second-Moment Inflation in Adaptive Optimization}
\label{subsec:optimizer_bias}

We now quantify how activation bias affects adaptive optimizers.
While gradient outputs exhibit no statistically significant mean
(Section~2.1), the activation tensor does.
Because per-token gradients are outer products
\[
g_i = \delta_i x_i^\top,
\]
any dense mean component in $x_i$ directly propagates into gradient
statistics.

\paragraph{Adam second-moment structure.}
Adam and AdamW maintain an elementwise second-moment estimate
\[
v \approx \mathbb{E}[g_i \odot g_i].
\]
Using the identity
\[
v = \frac{1}{N}(\Delta^{\odot 2})^\top (X^{\odot 2}),
\]
we see that squared activation magnitudes directly influence the
scale of $v$.

If $x = \mu_x + \tilde{x}$, then
\[
g_i = \delta_i \mu_x^\top + \delta_i \tilde{x}_i^\top,
\]
and the first term produces a dense, coherent contribution to $v$.

\paragraph{Experimental protocol.}
For each layer and training stage we compute:
\begin{itemize}
\item $v_{\text{raw}}$ from original activations,
\item $v_{\text{mean}}$ after mean removal,
\item $v_{\text{mean+whiten}}$ after mean removal and top-$1\%$ spike whitening.
\end{itemize}
We report extreme-tail statistics ($p99.9$ and $\max$), which dominate
learning-rate suppression in Adam.

\begin{table}[t]
\centering
\caption{Extreme-tail statistics of Adam second moment $v$.}
\label{tab:adam_v_bias}
\begin{tabular}{lcccc}
\toprule
Layer & Step & Method & $p99.9(v)$ & $\max(v)$ \\
\midrule
Layer 3 & 10k & Raw & $1.60\times10^{-10}$ & $8.38\times10^{-10}$ \\
 & & Mean-only & $3.44\times10^{-11}$ & $1.91\times10^{-10}$ \\
 & & Mean+whiten & $1.72\times10^{-11}$ & $8.09\times10^{-11}$ \\
\midrule
Layer 3 & 180k & Raw & $1.00\times10^{-10}$ & $7.37\times10^{-9}$ \\
 & & Mean-only & $2.50\times10^{-11}$ & $1.53\times10^{-10}$ \\
 & & Mean+whiten & $2.44\times10^{-11}$ & $1.54\times10^{-10}$ \\
\midrule
Layer 21 & 10k & Raw & $1.67\times10^{-11}$ & $3.43\times10^{-11}$ \\
 & & Mean-only & $1.59\times10^{-11}$ & $2.58\times10^{-11}$ \\
 & & Mean+whiten & $1.48\times10^{-11}$ & $2.33\times10^{-11}$ \\
\midrule
Layer 21 & 180k & Raw & $9.23\times10^{-11}$ & $1.10\times10^{-9}$ \\
 & & Mean-only & $2.21\times10^{-11}$ & $2.69\times10^{-10}$ \\
 & & Mean+whiten & $2.09\times10^{-11}$ & $1.95\times10^{-10}$ \\
\bottomrule
\end{tabular}
\end{table}

\paragraph{Key observations.}

\begin{enumerate}
\item \textbf{Mean removal dramatically reduces extreme $v$ outliers.}
In early layers and early training (Layer 3, 10k), the $p99.9$
statistic decreases by more than $4\times$, and the maximum by more
than $4\times$.
In some later-stage cases, the reduction in $\max(v)$ approaches
$50\times$.

\item \textbf{Median values change minimally.}
Across all runs, the median of $v$ changes by less than 1\%.
This indicates that bias primarily inflates a small subset of
extreme coordinates rather than the bulk of parameters.

\item \textbf{Whitening provides additional but smaller gains.}
In early training stages, spike whitening further reduces extreme
values.
In later stages, most of the pathological inflation is already
removed by mean subtraction.
\end{enumerate}

\paragraph{Interpretation.}

These results demonstrate that activation bias is a dominant driver
of extreme second-moment inflation in adaptive optimizers.
Even though gradient outputs themselves are nearly zero-mean,
the dense activation bias propagates multiplicatively through
outer-product gradients and produces heavy-tailed $v$ statistics.

Removing the activation mean significantly stabilizes second-moment
estimates and therefore mitigates excessive effective learning-rate
suppression in Adam-style optimizers.

This provides a direct mechanistic link between activation bias
and optimization pathologies.


\subsection{Bias-Induced Dynamic Range Distortion in Low-Bit Quantization}
\label{subsec:quant_bias}

Low-bit quantization is fundamentally sensitive to dynamic range.
Under blockwise quantization, each block of size $B$ uses a shared
scale parameter proportional to
\[
\alpha = \max_{j \in \text{block}} |x_j|.
\]
If a small subset of entries exhibits unusually large magnitude,
they determine $\alpha$, thereby compressing the effective resolution
allocated to the majority of entries.

Given the decomposition
\[
x = \mu_x + x_{\text{spike}} + x_{\text{tail}},
\]
we now empirically evaluate how much of the dynamic range distortion
is attributable to activation bias.

\paragraph{Experimental protocol.}
For each layer and training stage we compute:

\begin{itemize}
\item Entrywise magnitude statistics ($p99$, $p99.9$, $\max$),
\item Blockwise $p99.9$ of block-maximum magnitudes (block size 64),
\item Reduction ratios after mean removal,
\item Additional reductions after spike whitening.
\end{itemize}

\begin{table}[t]
\centering
\caption{Dynamic range statistics before and after bias removal.}
\label{tab:quant_bias}
\begin{tabular}{lcccc}
\toprule
Layer & Step & Method & Entry $p99.9$ & Block $p99.9$ \\
\midrule
Layer 3 & 10k & Raw & $6.24$ & $10.40$ \\
 & & Mean-only & $3.09$ & $5.27$ \\
 & & Mean+whiten & $2.34$ & $3.22$ \\
\midrule
Layer 3 & 180k & Raw & $2.29$ & $3.42$ \\
 & & Mean-only & $1.16$ & $1.71$ \\
 & & Mean+whiten & $1.06$ & $1.61$ \\
\midrule
Layer 21 & 10k & Raw & $3.52$ & $4.60$ \\
 & & Mean-only & $3.46$ & $4.52$ \\
 & & Mean+whiten & $4.00$ & $4.00$ \\
\midrule
Layer 21 & 180k & Raw & $4.17$ & $6.48$ \\
 & & Mean-only & $2.31$ & $4.20$ \\
 & & Mean+whiten & $2.98$ & $2.98$ \\
\bottomrule
\end{tabular}
\end{table}

\paragraph{Key observations.}

\begin{enumerate}
\item \textbf{Mean removal consistently reduces extreme magnitudes.}
In early-layer early-training regimes (Layer 3, 10k),
entrywise $p99.9$ reduces by approximately $2\times$,
and blockwise $p99.9$ reduces by approximately $2\times$.
This indicates that a substantial portion of dynamic range inflation
originates from the mean component.

\item \textbf{Later layers exhibit smaller bias sensitivity.}
Layer 21 at 10k shows minimal reduction after mean removal,
consistent with the weaker bias ratio observed in Section~2.1.

\item \textbf{Spike whitening further smooths residual anisotropy.}
In some regimes, spike whitening provides additional reduction
beyond mean removal, but the primary distortion mechanism remains
bias-dominated.
\end{enumerate}

\paragraph{Interpretation.}

These results demonstrate that activation bias directly inflates
the blockwise dynamic range that governs low-bit quantization.
Removing the bias reduces extreme magnitudes by approximately
$2\times$ in sensitive regimes, thereby restoring effective
resolution for long-tail semantic components.

Importantly, the majority of dynamic range distortion is attributable
to the mean component rather than the long-tail spectrum itself.

This provides a concrete mechanistic explanation for why mean removal
substantially improves training stability in FP4 and related
low-bit settings.
\section{A Formal Model of Bias, Spike Variance, and Tail Variability}
\label{sec:formal_model}

Section~\ref{sec:empirical} established three empirical facts:
(i) activations have a statistically significant mean ($R_x\gg 1$) while gradient
outputs are mean-free ($R_\delta\approx 1$),
(ii) extreme activation entries are dominated by the mean component,
and (iii) mean removal reduces but does not eliminate the spectral spike.
In this section we formalize these observations into a minimal model that
separates \emph{bias}, \emph{structured low-rank variance}, and a
\emph{high-dimensional tail}.
The model is intentionally simple: its purpose is to isolate the mechanisms
needed to explain the empirical phenomena and motivate later algorithmic
interventions.

\subsection{Token-wise representation model}
\label{subsec:model_token}

Let $x_i \in \mathbb{R}^H$ denote the representation (activation) of token $i$
at a fixed layer, and let
\[
X = \begin{bmatrix} x_1^\top \\ \vdots \\ x_N^\top \end{bmatrix}\in\mathbb{R}^{N\times H}
\]
be the stacked activation matrix for a batch of $N$ tokens.
We decompose each token representation as
\begin{equation}
\label{eq:x_decompose}
x_i = \mu + s_i,
\qquad
\mu \in \mathbb{R}^H,
\qquad
\mathbb{E}[s_i]=0,
\end{equation}
where $\mu$ is a shared bias component and $s_i$ captures token-dependent
variation.
Equivalently,
\[
X = M + S,
\qquad
M \triangleq \mathbf{1}\mu^\top,
\qquad
S \triangleq \begin{bmatrix} s_1^\top \\ \vdots \\ s_N^\top \end{bmatrix}.
\]
By construction, mean-centering removes the bias:
\[
X_c \triangleq X - M = S.
\]

\subsection{Spike--tail variance decomposition}
\label{subsec:model_spike_tail}

The empirical results indicate that $S$ (the centered activations) exhibits a
spike--tail spectrum.
We model this by assuming that the centered variation admits an orthogonal
decomposition into a low-rank structured component and a high-dimensional tail:
\begin{equation}
\label{eq:s_decompose}
s_i = s_i^{\text{spike}} + s_i^{\text{tail}},
\qquad
s_i^{\text{spike}} \in \mathcal{U},
\qquad
s_i^{\text{tail}} \in \mathcal{U}^\perp,
\end{equation}
for a low-dimensional subspace $\mathcal{U}\subset\mathbb{R}^H$ with
$\dim(\mathcal{U})=r \ll H$.
Let $P_{\mathcal{U}}$ and $P_{\mathcal{U}^\perp}$ denote the corresponding
orthogonal projectors. Then
\[
s_i^{\text{spike}} = P_{\mathcal{U}}s_i,
\qquad
s_i^{\text{tail}}  = P_{\mathcal{U}^\perp}s_i.
\]

In matrix form,
\[
S = S_{\text{spike}} + S_{\text{tail}},
\qquad
S_{\text{spike}} = SP_{\mathcal{U}},
\qquad
S_{\text{tail}} = SP_{\mathcal{U}^\perp}.
\]

\paragraph{Empirical correspondence.}
In practice, $\mathcal{U}$ is instantiated as the span of the top-$k$ right singular
vectors of the centered activation matrix $X_c$, where $k$ is chosen as a small
fraction of $H$ (e.g., $k=1\%$).
Section~\ref{subsec:spectral_role_bias} showed that mean removal reduces the
leading singular values but leaves a nontrivial residual low-rank variance
subspace, consistent with~\eqref{eq:s_decompose}.

\subsection{Elementwise outliers and the role of bias}
\label{subsec:model_outliers}

The model~\eqref{eq:x_decompose}--\eqref{eq:s_decompose} distinguishes between
two qualitatively different sources of large elementwise magnitudes:
\begin{itemize}
\item \textbf{Bias-driven outliers.} The bias term $\mu$ is dense and shared,
so the mean matrix $M=\mathbf{1}\mu^\top$ can contribute coherently to the
magnitude of many entries.
\item \textbf{Variance-driven extremes.} The random components
$s_i^{\text{spike}}$ and $s_i^{\text{tail}}$ fluctuate across tokens and can
produce large entries through variance, but do not contribute to the mean.
\end{itemize}
Section~\ref{subsec:outlier_mean_dominance} established empirically that, among
the largest $0.1\%$ entries of $X$, squared-magnitude contributions are dominated
by $M$, with smaller contributions from $S_{\text{spike}}$ and negligible
contributions from $S_{\text{tail}}$.

\subsection{Gradient decomposition under mean-biased activations}
\label{subsec:model_grad}

Let $\delta_i \in \mathbb{R}^{H'}$ denote the gradient output (backpropagated
signal) for token $i$ at a linear layer with weight matrix
$W\in\mathbb{R}^{H'\times H}$.
The per-token gradient is
\[
g_i = \delta_i x_i^\top \in \mathbb{R}^{H'\times H},
\]
and the batch gradient is
\[
G \triangleq \frac{1}{N}\sum_{i=1}^N g_i
= \mathbb{E}[\delta x^\top].
\]
Substituting~\eqref{eq:x_decompose} yields the identity
\begin{equation}
\label{eq:G_decompose}
G
=
\mathbb{E}[\delta]\mu^\top
+
\mathbb{E}[\delta s^\top].
\end{equation}
This separates a rank-one bias-driven term from a cross-covariance term.
Section~\ref{subsec:first_order_asymmetry} established empirically that
$\mathbb{E}[\delta]$ is noise-level, while $\mu$ is large.
Section~\ref{subsec:grad_spectrum_bias} further showed that centering $x$
reduces the spectral spike of $G$, consistent with removing the first term.

\subsection{Second-moment (Adam) statistics}
\label{subsec:model_adam}

Adaptive optimizers maintain elementwise second moments of gradients.
For the outer-product gradient $g_i=\delta_i x_i^\top$, the elementwise second
moment admits the exact batch identity
\begin{equation}
\label{eq:v_identity}
v \triangleq \mathbb{E}[g_i \odot g_i]
=
\frac{1}{N}(\Delta^{\odot 2})^\top (X^{\odot 2}),
\end{equation}
where $\Delta\in\mathbb{R}^{N\times H'}$ stacks $\delta_i$ and
$X\in\mathbb{R}^{N\times H}$ stacks $x_i$.
Equation~\eqref{eq:v_identity} makes explicit that activation magnitudes
directly influence second-moment inflation.
Section~\ref{subsec:optimizer_bias} showed empirically that mean removal reduces
the extreme tail of $v$ (up to $50\times$ in $\max(v)$ in some regimes),
establishing bias-driven second-moment inflation as a key mechanism.

\paragraph{Summary.}
The model isolates three components with distinct roles:
\begin{enumerate}
\item $\mu$ (bias): drives elementwise outliers and inflates second moments;
\item $\mathcal{U}$ (spike variance subspace): captures low-rank structured variance;
\item $\mathcal{U}^\perp$ (tail): carries high-dimensional token-dependent variability.
\end{enumerate}
The next section studies algorithmic consequences: mean subtraction removes
$\mu$ exactly with negligible cost, while lightweight spectral conditioning can
reduce residual spike variance when needed.
\section{Theoretical Consequences of Bias}
\label{sec:theory}

Section~\ref{sec:formal_model} introduced a structured decomposition of token
representations:
\[
x_i = \mu + s_i^{\text{spike}} + s_i^{\text{tail}},
\]
where $\mu$ is a shared bias component, $s_i^{\text{spike}}$ lies in a
low-dimensional subspace, and $s_i^{\text{tail}}$ occupies a high-dimensional
complement.
Section~\ref{sec:empirical} established empirically that activation bias
dominates extreme entries and significantly inflates optimizer statistics.
We now formalize these observations.

Throughout this section, we assume:
\begin{itemize}
\item $\mu$ is dense and fixed across tokens,
\item $s_i$ are zero-mean random vectors with finite second moment,
\item $s_i$ have sub-Gaussian tails (for concentration results).
\end{itemize}

The assumptions reflect the empirical statistics observed in Section~2.

\subsection{Bias Dominates Extreme Activation Entries}
\label{subsec:theory_outliers}

We begin by formalizing why bias dominates the largest-magnitude activation
entries.

\begin{theorem}[Bias dominance of extreme entries]
\label{thm:bias_outlier}
Let $x_i = \mu + s_i$ with $\mathbb{E}[s_i]=0$ and
$\|s_{ij}\|$ sub-Gaussian with variance proxy $\sigma^2$.
If $\max_j |\mu_j| \gg \sigma$,
then for sufficiently small $\varepsilon>0$, among the top
$\varepsilon$-fraction of entries (by $|x_{ij}|$),
\[
\frac{|M_{ij}|^2}{|x_{ij}|^2} \to 1
\quad \text{with high probability as } N\to\infty.
\]
\end{theorem}

\paragraph{Interpretation.}
Extreme activation entries are primarily caused by the coherent bias component.
Random fluctuations from $s_i$ produce heavy tails only with exponentially small
probability, whereas the bias contributes deterministically to every token.

\paragraph{Connection to Section~2.2.}
The empirical result that mean contribution exceeds $85\%$ of squared magnitude
for top-$0.1\%$ entries (Section~2.2) is a finite-sample manifestation of
Theorem~\ref{thm:bias_outlier}.

\subsection{Bias as a Structured Component of the Spectral Spike}
\label{subsec:theory_spike}

The activation matrix can be written as
\[
X = M + S,
\qquad
M = \mathbf{1}\mu^\top.
\]

\begin{theorem}[Bias contributes a rank-one spike]
\label{thm:bias_spike}
The matrix $M$ has rank one with singular value
\[
\sigma_1(M) = \sqrt{N}\|\mu\|_2.
\]
Moreover,
\[
\sigma_1(X) \ge \sigma_1(M) - \|S\|_2.
\]
If $\sqrt{N}\|\mu\|_2$ dominates $\|S\|_2$,
then the leading singular direction of $X$ aligns with $\mu$.
\end{theorem}

\paragraph{Interpretation.}
Bias produces a coherent low-rank component inside the spike subspace.
Mean removal eliminates this rank-one contribution exactly, but does not
eliminate structured variance in $S_{\text{spike}}$.
This explains why Section~2.3 observed partial but not complete spike reduction
after centering.

\subsection{Bias-Induced Gradient Structure}
\label{subsec:theory_grad}

Consider a linear layer with gradient
\[
G = \mathbb{E}[\delta x^\top].
\]
Substituting $x = \mu + s$ yields the exact identity
\begin{equation}
\label{eq:G_decomp_theory}
G =
\mathbb{E}[\delta]\mu^\top
+
\mathbb{E}[\delta s^\top].
\end{equation}

\begin{theorem}[Bias-driven gradient component]
\label{thm:grad_bias}
Even if $\mathbb{E}[\delta]=0$, the second-moment statistics of $G$
contain bias-dependent inflation terms:
\[
v = \mathbb{E}[g \odot g]
= \frac{1}{N}(\Delta^{\odot 2})^\top (X^{\odot 2}).
\]
If $X = M + S$, then
\[
X^{\odot 2}
=
M^{\odot 2}
+
2(M \odot S)
+
S^{\odot 2}.
\]
The first term contributes deterministic inflation proportional to $\mu_j^2$.
\end{theorem}

\paragraph{Interpretation.}
Bias increases extreme entries of $v$, shrinking effective learning rates in
Adam-type optimizers.
Mean removal reduces this inflation multiplicatively, consistent with the
empirical reductions observed in Section~2.4B.

\subsection{Dynamic-Range Inflation Under Bias}
\label{subsec:theory_quant}

Blockwise quantization uses scales proportional to
\[
\alpha = \max_{j \in \text{block}} |x_{ij}|.
\]

\begin{theorem}[Bias-driven dynamic range inflation]
\label{thm:quant_bias}
Under the model $x_i=\mu+s_i$,
if $|\mu_j|$ dominates typical fluctuations of $s_{ij}$,
then for blocks containing dimension $j$,
\[
\alpha \approx |\mu_j|
\]
with high probability.
\end{theorem}

\paragraph{Interpretation.}
Bias inflates quantization scale and compresses semantic resolution.
Mean removal restores dynamic range symmetry, explaining the empirical
improvements observed in Section~2.4C.

\paragraph{Summary of theoretical implications.}

The preceding results establish three mechanisms:

\begin{enumerate}
\item Bias dominates extreme activation entries.
\item Bias contributes a coherent low-rank spectral component.
\item Bias inflates second-moment and quantization statistics.
\end{enumerate}

These mechanisms collectively explain why mean removal is an effective
and computationally trivial intervention.
Detailed proofs of Theorems~\ref{thm:bias_outlier}--\ref{thm:quant_bias}
are provided in Appendix~A.


\section{Practical Implications: Bias as a Controllable Source of Instability}
\label{sec:implications}

The preceding sections established three key empirical and theoretical facts:

\begin{enumerate}
    \item Activation bias is statistically significant and structurally coherent,
    while gradient output mean remains at noise level.
    \item Extreme activation outliers are overwhelmingly dominated by the bias
    component rather than spike variance or tail variability.
    \item Bias inflates both blockwise quantization scale and adaptive optimizer
    second-moment statistics.
\end{enumerate}

These observations allow us to reinterpret several long-standing training
pathologies through a unified lens: many instabilities attributed to spectral
anisotropy are in fact primarily driven by activation bias.

We now discuss the implications for quantization and optimization.

\subsection{Quantization: Bias as Dynamic-Range Inflation}

Low-bit quantization operates under a fundamental constraint: each block of
features shares a common scaling factor. The effective resolution allocated to
each element is therefore inversely proportional to the blockwise scale.

From Appendix~\ref{app:quant_bias}, if a block contains coordinates with large
$|\mu_j|$, the scale concentrates around
\[
\alpha(x;B) \approx \max_{j \in B} |\mu_j|,
\]
rather than the intrinsic stochastic variability level
$O(\sigma \sqrt{\log b})$.

This inflation has two consequences:

\begin{itemize}
    \item The quantization step size increases proportionally.
    \item The long-tail semantic component $s_j$ is compressed relative to
    its intrinsic scale.
\end{itemize}

Empirically, mean removal reduces blockwise $p99.9$ magnitudes by
approximately a factor of two in early layers, consistent with the
high-probability bounds derived in Appendix~\ref{app:quant_bias}.
The resulting reduction in quantization step size directly improves effective
tail signal-to-noise ratio (SNR).

Importantly, this mechanism does not require full spike whitening.
Bias removal alone addresses the dominant source of dynamic-range inflation.
Spike whitening may provide further improvements when spike variance remains
substantial, but it is not necessary to eliminate the principal outlier driver.

\subsection{Optimization: Bias as Second-Moment Inflation}

Adaptive optimizers such as Adam maintain elementwise second-moment estimates
\[
v_{ab} = \mathbb{E}[g_{ab}^2],
\qquad
g_{ab} = \delta_a x_b.
\]

Appendix~\ref{app:gradient_proof} shows that
\[
v_{ab}
=
\mu_b^2 \mathbb{E}[\delta_a^2]
+
\mathbb{E}[\delta_a^2 s_b^2],
\]
under mild independence assumptions.

When $\mu_b^2$ dominates $\mathbb{E}[s_b^2]$, the second moment becomes
bias-driven. This inflates $v_{ab}$ for specific coordinates, thereby
suppressing their effective learning rates via the update
\[
\Delta W_{ab} \propto \frac{g_{ab}}{\sqrt{v_{ab}}}.
\]

Empirically, mean removal reduces the extreme tail of $v$ by factors ranging
from $4\times$ to nearly $50\times$ in early layers, while leaving median
values largely unchanged. This confirms that bias primarily affects a small
subset of coordinates that nonetheless dominate adaptive learning-rate
suppression.

Unlike gradient-output mean (which is statistically noise-level), activation
bias propagates multiplicatively through the outer-product structure
$g_i = \delta_i x_i^\top$ and therefore directly controls second-moment
inflation.

\subsection{When Is Mean Removal Sufficient?}

The experiments suggest a layered interpretation:

\begin{itemize}
    \item \textbf{Early layers / early training:} bias dominates both
    outliers and second-moment inflation.
    Mean removal alone resolves most instability.
    
    \item \textbf{Later layers / late training:} spike variance persists even
    after bias removal.
    Additional spectral whitening may further flatten the spectrum.
\end{itemize}

Thus, bias removal can be viewed as the minimal structural correction.
Spike whitening becomes an incremental refinement rather than the primary fix.

\subsection{Relation to Existing Spectral Methods}

Full SVD-based spike removal eliminates both bias and spike variance.
Our analysis clarifies that a substantial portion of the benefit previously
attributed to spike suppression may arise from removing the coherent mean
component.

This observation has two practical consequences:

\begin{enumerate}
    \item In low-bit training, mean removal can replace expensive SVD
    decompositions when bias dominates dynamic-range inflation.
    
    \item In optimizer design, centering activations may mitigate second-moment
    inflation without requiring explicit spike tracking.
\end{enumerate}

These implications motivate a re-examination of spectral regularization
techniques: in many regimes, controlling bias may be sufficient to stabilize
training.

\subsection{Summary}

The empirical and theoretical analyses converge on a consistent picture:

\begin{itemize}
    \item Bias is the primary driver of extreme activation outliers.
    \item Bias propagates multiplicatively into gradient second moments.
    \item Removing bias reduces both quantization instability and optimizer
    suppression.
    \item Spike variance remains but plays a secondary role relative to bias.
\end{itemize}

These results reposition mean removal from a heuristic preprocessing step to a
structurally justified intervention grounded in activation geometry.
\section{Bias Mitigation Improves Multiple Downstream Training Regimes}
\label{sec:downstream}

Sections~2--4 established that activation bias is:
\begin{itemize}
    \item statistically significant in activations,
    \item responsible for the majority of extreme activation outliers,
    \item a dominant contributor to second-moment inflation,
    \item and a structured component inside the spectral spike.
\end{itemize}

We now evaluate whether mitigating bias improves training across
multiple regimes.
The goal is not to solve a single isolated problem,
but to demonstrate that bias control has broad consequences
for large language model training.

All experiments use a $1$B-parameter Transformer trained for $100$B tokens.

\subsection{Low-Bit FP4 Training}

We first evaluate FP4 blockwise quantization.
The methods compared are:

\begin{enumerate}
    \item Raw FP4 (no conditioning)
    \item Mean-only removal
    \item Mean + spike whitening (top $1\%$)
    \item Metis (SVD-based spike suppression)
    \item FP16 baseline
\end{enumerate}

\paragraph{Training stability.}

Raw FP4 diverges early due to dynamic-range inflation.
Mean removal alone stabilizes training.
Mean + whitening further improves convergence smoothness.

Figure~\ref{fig:fp4_curves} shows training loss over $100$B tokens.

\begin{figure}[t]
\centering
\includegraphics[width=0.9\linewidth]{figures/fp4_training_curves.png}
\caption{FP4 training curves.
Mean removal stabilizes training relative to raw FP4.
Mean + whitening approaches SVD-based Metis performance.}
\label{fig:fp4_curves}
\end{figure}

\paragraph{Interpretation.}

The improvement is consistent with Section~2.2:
outliers are predominantly bias-driven.
Removing bias reduces blockwise scaling inflation.
Whitening refines the remaining spike variance.

Importantly, the majority of stabilization comes from bias removal,
not from full SVD suppression.

\subsection{Adaptive Optimization (AdamW)}

We next evaluate bias mitigation under full precision
using AdamW-style optimization.

Compared methods:

\begin{enumerate}
    \item AdamW baseline
    \item AdamW + activation mean removal
    \item Spectra (spike-aware optimizer)
\end{enumerate}

\paragraph{Training behavior.}

Figure~\ref{fig:opt_curves} shows convergence behavior.

\begin{figure}[t]
\centering
\includegraphics[width=0.9\linewidth]{figures/optimizer_training_curves.png}
\caption{Optimizer comparison.
Mean-centered AdamW improves stability and narrows the gap to Spectra.}
\label{fig:opt_curves}
\end{figure}

\paragraph{Interpretation.}

From Section~2.4,
bias inflates the extreme tail of the second-moment estimate $v$.
Mean removal reduces this inflation,
thereby increasing effective learning rates
without altering semantic covariance structure.

The improvement is strongest in early layers and early training,
where activation bias is largest.

\subsection{Cross-Regime Consistency}

The same structural phenomenon explains improvements in both settings:

\begin{itemize}
    \item In quantization, bias inflates dynamic range.
    \item In optimization, bias inflates second-moment statistics.
\end{itemize}

The mitigation is identical:
remove the coherent mean component.

This cross-regime consistency supports the central claim:
bias is a shared structural driver of multiple training pathologies.

\subsection{Summary of Downstream Effects}

\begin{table}[t]
\centering
\caption{Summary of bias mitigation effects.}
\label{tab:downstream_summary}
\begin{tabular}{lcc}
\toprule
Setting & Raw & Mean-Removed \\
\midrule
FP4 stability & unstable & stable \\
Block dynamic range & inflated & reduced \\
Adam second moment tail & heavy & reduced \\
Effective LR allocation & suppressed & improved \\
\bottomrule
\end{tabular}
\end{table}

Across both regimes,
mean removal produces consistent stabilization effects.

Spike whitening provides additional refinement,
but bias control is the dominant intervention.
\section{Conclusion and Discussion}
\label{sec:conclusion}

This work re-examines anisotropy in large language model training
through the lens of activation bias.
While prior analyses focused on spectral spike--tail structure,
we show that a coherent mean component in activations
is the primary driver of extreme elementwise outliers
and second-moment inflation.

Empirically, we demonstrated that:

\begin{itemize}
    \item Activation means are statistically significant across layers and
    training stages, while gradient-output means remain at noise level.
    \item Extreme activation outliers are overwhelmingly dominated by
    the mean component rather than spike variance or semantic tail.
    \item Bias propagates multiplicatively into gradient outer products,
    inflating both blockwise quantization scale and adaptive optimizer
    second moments.
\end{itemize}

These findings clarify that many training instabilities attributed to
spectral anisotropy are, in practice, bias-driven.

\paragraph{Curse.}
Bias introduces dynamic-range inflation in low-bit quantization,
suppresses effective learning rates in adaptive optimizers,
and amplifies extreme coordinate statistics.
Its coherent nature causes disproportionate influence on
blockwise scaling and second-moment estimates.

\paragraph{Blessing.}
Because bias is structured, low-rank, and globally coherent,
it admits an extremely simple mitigation strategy:
mean removal.
Unlike full SVD-based spike suppression,
mean removal is hardware-friendly, inexpensive,
and conservative — it removes the dominant instability
without disturbing semantic covariance structure.

Across large-scale experiments ($1$B parameters, $100$B tokens),
bias mitigation improved both FP4 quantization stability
and adaptive optimizer behavior.
Spike whitening provided further refinement,
but the majority of stabilization arose from bias removal alone.

\subsection{Implications}

The results suggest a revised perspective on LLM training geometry:

\begin{enumerate}
    \item Spectral spike structure is not monolithic.
    A significant portion of spike energy corresponds to bias.
    \item Controlling bias is a minimal structural correction
    that addresses multiple downstream instabilities simultaneously.
    \item Expensive spectral decompositions may be solving
    a simpler underlying structural issue.
\end{enumerate}

This reframing opens new design directions
for quantization schemes, optimizers,
and potentially even architectural normalization strategies.

\subsection{Limitations and Open Questions}

Several important questions remain open:

\begin{itemize}
    \item \textbf{Origin of activation bias.}
    Why do forward activations consistently develop large coherent means?
    Is this purely data-driven, architecture-induced,
    or an emergent property of residual accumulation?
    
    \item \textbf{Interaction with normalization layers.}
    How does bias behave under different normalization schemes
    (pre-norm vs post-norm, RMSNorm vs LayerNorm)?
    
    \item \textbf{Scaling behavior.}
    Does bias magnitude scale with model size or training duration?
    Preliminary evidence suggests early-layer dominance,
    but a systematic scaling law analysis is warranted.
    
    \item \textbf{Beyond mean removal.}
    Can lightweight structured corrections approximate full whitening
    without expensive decompositions?
\end{itemize}

Addressing these questions will further clarify the geometric
role of bias in representation learning.

\subsection{Final Perspective}

Bias is neither purely pathological nor purely benign.
It simultaneously destabilizes training
and provides a coherent structural handle for mitigation.

Recognizing this duality
allows us to move beyond viewing anisotropy
as an opaque spectral phenomenon,
toward understanding and controlling
the specific geometric mechanisms that drive it.

In this sense, bias embodies both
the curse and the blessing
of large language model training.

\newpage

\bibliographystyle{../../doc/iclr-template/master/iclr2026/iclr2026_conference}
\bibliography{references}


\newpage

\appendix
\appendix
\section{Proof of Mean-Dominance for Extreme Entries}
\label{app:outlier_bias_proof}

This appendix proves a formal version of the empirical observation in
Section~\ref{subsec:outlier_mean_dominance}:
among the largest-magnitude entries of the activation matrix, the broadcasted
mean component dominates.

\subsection{Setup and notation}

Let $x_i \in \mathbb{R}^H$ denote token activations and let $X\in\mathbb{R}^{N\times H}$
stack them row-wise. We model each entry as
\begin{equation}
\label{eq:app_model_entry}
x_{ij} = \mu_j + s_{ij},
\end{equation}
where $\mu\in\mathbb{R}^H$ is deterministic (the coherent mean) and
$\{s_{ij}\}$ are random fluctuations with $\mathbb{E}[s_{ij}]=0$.

Define the mean matrix $M\triangleq \mathbf{1}\mu^\top$ and the centered matrix
$S\triangleq X-M$ so that $S_{ij}=s_{ij}$.

For any $\varepsilon\in(0,1)$, define the (population) outlier set as the indices
of the largest $\varepsilon NH$ entries by magnitude:
\begin{equation}
\label{eq:app_outlier_set}
\Omega_\varepsilon
\;\triangleq\;
\left\{(i,j): |x_{ij}| \text{ is among the top } \varepsilon NH \text{ values of } \{|x_{i'j'}|\}_{i',j'} \right\}.
\end{equation}
We will show that under mild conditions, entries in $\Omega_\varepsilon$ have
$s_{ij}$ small relative to $\mu_j$, hence the mean dominates.

\subsection{Assumptions}

\begin{assumption}[Sub-Gaussian fluctuations]
\label{assump:subg}
The entries $\{s_{ij}\}$ are independent and sub-Gaussian with parameter $\sigma^2$:
for all $t\ge 0$,
\[
\mathbb{P}\big(|s_{ij}| \ge t\big) \le 2\exp\!\left(-\frac{t^2}{2\sigma^2}\right).
\]
\end{assumption}

\begin{assumption}[Mean-gap and mass near the maximum]
\label{assump:gap}
Let $m \triangleq \max_{1\le j\le H} |\mu_j|$.
Fix $\gamma\in(0,1)$ and define the ``high-mean'' coordinate set
\[
\mathcal{J}_\gamma \triangleq \{j: |\mu_j|\ge (1-\gamma)m\},
\qquad
p_\gamma \triangleq |\mathcal{J}_\gamma|/H.
\]
Assume $p_\gamma>0$ and that for $j\notin\mathcal{J}_\gamma$, we have the strict gap
$|\mu_j|\le (1-\gamma)m$.
\end{assumption}

\begin{assumption}[Spike-dominance regime]
\label{assump:regime}
The mean amplitude satisfies
\[
m \;\ge\; \frac{4\sigma}{\gamma}\sqrt{2\log\!\left(\frac{4}{\varepsilon p_\gamma}\right)}.
\]
\end{assumption}

Assumption~\ref{assump:gap} captures the empirically observed phenomenon that a
nontrivial fraction of coordinates have mean values close to the maximum mean.
Assumption~\ref{assump:regime} places us in a regime where the mean dominates
rare noise events at the outlier quantile.

\subsection{Main theorem}

\begin{theorem}[Mean dominates the top-$\varepsilon$ fraction of entries]
\label{thm:mean_dominates_outliers}
Under Assumptions~\ref{assump:subg}--\ref{assump:regime}, with probability at least
\[
1 - 2\exp\!\left(-\frac{\varepsilon p_\gamma NH}{32}\right) - 2\varepsilon NH \exp\!\left(-\frac{\gamma^2 m^2}{32\sigma^2}\right),
\]
the following two statements hold:

\begin{enumerate}
\item (\emph{Outliers occur on high-mean coordinates.})
Every outlier index $(i,j)\in \Omega_\varepsilon$ satisfies $j\in \mathcal{J}_\gamma$
and
\[
|s_{ij}| \le \frac{\gamma m}{4}.
\]

\item (\emph{Per-outlier mean dominance.})
For every $(i,j)\in \Omega_\varepsilon$,
the squared-magnitude share of the mean satisfies
\[
\frac{M_{ij}^2}{M_{ij}^2 + S_{ij}^2}
\;\ge\;
1 - \left(\frac{\gamma}{4(1-\gamma)}\right)^2.
\]
\end{enumerate}
\end{theorem}

\subsection{Proof}

We prove two counting lemmas: (i) there are enough ``clean'' entries on high-mean
coordinates to fill the outlier set, and (ii) low-mean coordinates cannot produce
enough entries above the outlier threshold without requiring rare large noise.

\begin{lemma}[Many high-mean entries have small noise]
\label{lem:many_clean_highmean}
Let $t \triangleq \gamma m/4$.
Then with probability at least $1-2\exp(-\varepsilon p_\gamma NH/32)$, the number of
indices $(i,j)$ with $j\in\mathcal{J}_\gamma$ and $|s_{ij}|\le t$ is at least
\[
\frac{1}{2}p_\gamma NH.
\]
\end{lemma}

\begin{proof}
For each $(i,j)$ with $j\in\mathcal{J}_\gamma$, define the indicator
$Z_{ij} \triangleq \mathbf{1}\{|s_{ij}|\le t\}$.
By Assumption~\ref{assump:subg},
\[
\mathbb{P}(Z_{ij}=1) = 1-\mathbb{P}(|s_{ij}|>t)
\;\ge\;
1-2\exp\!\left(-\frac{t^2}{2\sigma^2}\right).
\]
Under Assumption~\ref{assump:regime}, $t=\gamma m/4$ implies
$2\exp(-t^2/(2\sigma^2)) \le \varepsilon p_\gamma/2$,
hence $\mathbb{P}(Z_{ij}=1)\ge 1-\varepsilon p_\gamma/2\ge 1/2$.
The total count is
\[
Z \triangleq \sum_{j\in\mathcal{J}_\gamma}\sum_{i=1}^N Z_{ij},
\qquad
\mathbb{E}[Z] \ge \frac{1}{2} p_\gamma NH.
\]
Since the $Z_{ij}$ are independent Bernoulli variables, a Chernoff bound yields
\[
\mathbb{P}\!\left(Z \le \frac{1}{2}\mathbb{E}[Z]\right)
\le
\exp\!\left(-\frac{\mathbb{E}[Z]}{8}\right)
\le
\exp\!\left(-\frac{p_\gamma NH}{32}\right)
\le
\exp\!\left(-\frac{\varepsilon p_\gamma NH}{32}\right),
\]
which proves the claim (up to a constant-factor tightening; we use the displayed bound).
\end{proof}

\begin{lemma}[Low-mean coordinates rarely exceed the high-mean threshold]
\label{lem:lowmean_rare}
Let $t=\gamma m/4$ and define the deterministic threshold
\[
\tau \triangleq (1-\gamma)m - t.
\]
For any $j\notin\mathcal{J}_\gamma$ and any $i$,
\[
\mathbb{P}\big(|x_{ij}|\ge \tau\big)
\le
2\exp\!\left(-\frac{t^2}{2\sigma^2}\right).
\]
Consequently, with probability at least $1-2\varepsilon NH\exp(-t^2/(2\sigma^2))$,
the number of indices $(i,j)$ with $j\notin\mathcal{J}_\gamma$ and $|x_{ij}|\ge\tau$
is at most $\varepsilon NH$.
\end{lemma}

\begin{proof}
Fix $j\notin\mathcal{J}_\gamma$, so $|\mu_j|\le (1-\gamma)m$ by Assumption~\ref{assump:gap}.
If $|x_{ij}|=|\mu_j+s_{ij}|\ge \tau=(1-\gamma)m-t$, then necessarily
$|s_{ij}|\ge t$ (otherwise $|\mu_j+s_{ij}|\le |\mu_j|+|s_{ij}|\le (1-\gamma)m+t < \tau$).
Thus
\[
\mathbb{P}(|x_{ij}|\ge\tau)\le \mathbb{P}(|s_{ij}|\ge t)
\le 2\exp\!\left(-\frac{t^2}{2\sigma^2}\right).
\]
The second statement follows by a union bound over all such indices and Markov/Chernoff-style counting; we use a simple union bound:
\[
\mathbb{P}\left(\exists\,\varepsilon NH \text{ distinct indices with } j\notin\mathcal{J}_\gamma,\,|x_{ij}|\ge\tau\right)
\le
2\varepsilon NH\exp\!\left(-\frac{t^2}{2\sigma^2}\right).
\]
\end{proof}

\begin{proof}[Proof of Theorem~\ref{thm:mean_dominates_outliers}]
Let $t=\gamma m/4$ and $\tau=(1-\gamma)m-t$ as above.

By Lemma~\ref{lem:many_clean_highmean}, with high probability there exist at least
$\frac{1}{2}p_\gamma NH$ indices $(i,j)$ with $j\in\mathcal{J}_\gamma$ and
$|s_{ij}|\le t$.
For each such index, we have
\[
|x_{ij}| = |\mu_j+s_{ij}|
\ge |\mu_j|-|s_{ij}|
\ge (1-\gamma)m - t
= \tau.
\]
Therefore, there are at least $\frac{1}{2}p_\gamma NH$ entries with magnitude at least $\tau$.

Under Assumption~\ref{assump:regime}, $p_\gamma/2 \ge \varepsilon$ (equivalently, $p_\gamma \ge 2\varepsilon$),
so these entries alone exceed the required $\varepsilon NH$ outliers.
Hence the outlier set $\Omega_\varepsilon$ is contained within the set of indices
with $|x_{ij}|\ge\tau$.

Now apply Lemma~\ref{lem:lowmean_rare}:
with high probability, at most $\varepsilon NH$ entries from low-mean coordinates
($j\notin\mathcal{J}_\gamma$) can reach magnitude $\tau$.
Since $\Omega_\varepsilon$ consists of the top $\varepsilon NH$ entries,
it follows that, with the stated probability, every $(i,j)\in\Omega_\varepsilon$
must satisfy $j\in\mathcal{J}_\gamma$ and $|s_{ij}|\le t$.
This proves Item (1).

For Item (2), for $(i,j)\in\Omega_\varepsilon$ we have $|\mu_j|\ge (1-\gamma)m$ and
$|s_{ij}|\le t=\gamma m/4$, hence
\[
\frac{M_{ij}^2}{M_{ij}^2 + S_{ij}^2}
=
\frac{\mu_j^2}{\mu_j^2 + s_{ij}^2}
\ge
\frac{(1-\gamma)^2m^2}{(1-\gamma)^2m^2 + (\gamma m/4)^2}
=
1 - \left(\frac{\gamma}{4(1-\gamma)}\right)^2,
\]
as required.
\end{proof}

\paragraph{Remark.}
The theorem formalizes a qualitative mechanism: when a nontrivial fraction of coordinates
carry large coherent means and the noise is sub-Gaussian, the largest-magnitude entries
come primarily from those mean-heavy coordinates, so mean subtraction removes the dominant
source of outliers. The residual spike--tail structure then reflects variance anisotropy,
consistent with Section~\ref{subsec:spectral_role_bias}.
\section{Gradient Decomposition and Second-Moment Inflation}
\label{app:gradient_proof}

This appendix formalizes how activation bias propagates into gradients and
inflates adaptive optimizer statistics.

\subsection{Notation and setup}

Consider a linear layer with weight matrix $W \in \mathbb{R}^{d \times h}$,
input activation $x_i \in \mathbb{R}^{h}$,
and output gradient $\delta_i \in \mathbb{R}^{d}$.

The per-token gradient is
\[
g_i = \delta_i x_i^\top \in \mathbb{R}^{d \times h}.
\]

Let
\[
G = \frac{1}{N} \sum_{i=1}^N g_i
\]
denote the batch gradient.

Assume activations decompose as
\[
x_i = \mu + s_i,
\qquad
\mathbb{E}[s_i] = 0.
\]

No assumption is made about $\mathbb{E}[\delta_i]$ at this stage.

\subsection{Exact gradient decomposition}

\begin{lemma}[Gradient decomposition identity]
\label{lem:grad_decomp}
The batch gradient satisfies
\[
G
=
\mathbb{E}[\delta]\mu^\top
+
\mathbb{E}[\delta s^\top].
\]
\end{lemma}

\begin{proof}
Substitute $x_i = \mu + s_i$ into
\[
G = \frac{1}{N}\sum_{i=1}^N \delta_i x_i^\top
=
\frac{1}{N}\sum_{i=1}^N \delta_i(\mu^\top + s_i^\top).
\]
Distribute:
\[
G
=
\left(\frac{1}{N}\sum_{i=1}^N \delta_i\right)\mu^\top
+
\frac{1}{N}\sum_{i=1}^N \delta_i s_i^\top.
\]
Taking expectation over batch sampling yields the stated identity.
\end{proof}

\paragraph{Interpretation.}
Even if $\mathbb{E}[\delta]=0$, the second term
$\mathbb{E}[\delta s^\top]$ remains.
However, the first term is a rank-one matrix aligned with $\mu$ and can
produce coherent gradient structure whenever $\mathbb{E}[\delta]\neq 0$.

\subsection{Second-moment structure of gradients}

Adaptive optimizers such as Adam maintain an elementwise second-moment estimate:
\[
v_{ab} = \mathbb{E}[g_{i,ab}^2].
\]

Since
\[
g_{i,ab} = \delta_{i,a} x_{i,b},
\]
we have the exact identity
\[
v_{ab}
=
\mathbb{E}[\delta_{a}^2 x_b^2].
\]

Using $x_b = \mu_b + s_b$, we expand:

\begin{align}
x_b^2
&=
(\mu_b + s_b)^2 \\
&=
\mu_b^2
+
2\mu_b s_b
+
s_b^2.
\end{align}

Therefore,

\begin{equation}
\label{eq:v_expand}
v_{ab}
=
\mu_b^2 \mathbb{E}[\delta_a^2]
+
2\mu_b \mathbb{E}[\delta_a^2 s_b]
+
\mathbb{E}[\delta_a^2 s_b^2].
\end{equation}

\subsection{Bias-induced inflation term}

If $s_b$ is zero-mean and independent of $\delta_a^2$,
then $\mathbb{E}[\delta_a^2 s_b] = 0$, and we obtain:

\begin{equation}
\label{eq:v_clean}
v_{ab}
=
\mu_b^2 \mathbb{E}[\delta_a^2]
+
\mathbb{E}[\delta_a^2 s_b^2].
\end{equation}

The first term is deterministic and shared across tokens.

\begin{theorem}[Bias-driven second-moment inflation]
\label{thm:second_moment_inflation}
Suppose
\[
\mu_b^2 \gg \mathbb{E}[s_b^2]
\]
for a subset of coordinates $b$.
Then for those coordinates,
\[
v_{ab}
\approx
\mu_b^2 \mathbb{E}[\delta_a^2],
\]
and hence the second moment is dominated by the bias term.
\end{theorem}

\begin{proof}
From \eqref{eq:v_clean},
\[
v_{ab}
=
\mu_b^2 \mathbb{E}[\delta_a^2]
+
\mathbb{E}[\delta_a^2 s_b^2].
\]
If $\mu_b^2 \gg \mathbb{E}[s_b^2]$ and $\delta_a^2$ has finite variance,
then
\[
\mu_b^2 \mathbb{E}[\delta_a^2]
\gg
\mathbb{E}[\delta_a^2 s_b^2].
\]
Therefore the first term dominates.
\end{proof}

\paragraph{Consequences.}
Large $\mu_b$ causes certain coordinates to have inflated $v_{ab}$ values.
Since Adam uses effective learning rate proportional to $1/\sqrt{v_{ab}}$,
these coordinates experience suppressed learning rates.

\subsection{Effect of mean removal}

After centering activations,
\[
x'_i = x_i - \mu = s_i.
\]

The new second moment becomes

\[
v'_{ab}
=
\mathbb{E}[\delta_a^2 s_b^2].
\]

\begin{theorem}[Second-moment reduction under mean removal]
\label{thm:v_reduction}
For coordinates where bias dominates,
\[
\frac{v'_{ab}}{v_{ab}}
\approx
\frac{\mathbb{E}[s_b^2]}{\mu_b^2}
\ll 1.
\]
\end{theorem}

\begin{proof}
Divide \eqref{eq:v_clean}:
\[
\frac{v'_{ab}}{v_{ab}}
=
\frac{\mathbb{E}[\delta_a^2 s_b^2]}
{\mu_b^2 \mathbb{E}[\delta_a^2] + \mathbb{E}[\delta_a^2 s_b^2]}.
\]
If $\mu_b^2 \gg \mathbb{E}[s_b^2]$, denominator dominated by first term,
so
\[
\frac{v'_{ab}}{v_{ab}}
\approx
\frac{\mathbb{E}[s_b^2]}
{\mu_b^2}
\ll 1.
\]
\end{proof}

\paragraph{Interpretation.}
Mean removal reduces extreme $v_{ab}$ entries multiplicatively.
This matches empirical reductions of $4\times$ to $50\times$ observed in
Section~2.4B.

\subsection{Summary}

We have shown:

\begin{enumerate}
\item Bias induces a rank-one component in gradients.
\item Bias contributes a deterministic inflation term to second-moment statistics.
\item Mean removal eliminates this term exactly.
\item The reduction factor scales with $\mu_b^2 / \mathbb{E}[s_b^2}$.
\end{enumerate}

This provides a precise mechanism linking activation bias to optimizer
instability.
\section{Bias-Driven Dynamic-Range Inflation in Block Quantization}
\label{app:quant_bias}

This appendix formalizes how activation bias inflates blockwise quantization
scales and thereby degrades the effective resolution available to long-tail
components.

\subsection{Setup: block quantization model}

Consider a vector $x\in\mathbb{R}^H$ that will be quantized using blockwise
scaling.
Let $B\subset\{1,\dots,H\}$ denote a block of coordinates with $|B|=b$.
A standard block quantizer chooses a scale proportional to
\begin{equation}
\label{eq:block_scale_max}
\alpha(x;B) \triangleq \max_{j\in B} |x_j|.
\end{equation}
The quantization step size is then
\[
\Delta(x;B) = \alpha(x;B) / Q,
\]
where $Q$ is the maximum representable integer level (e.g.\ $Q=7$ for signed
4-bit symmetric quantization).

We study how $\alpha(x;B)$ behaves under the decomposition
\begin{equation}
\label{eq:x_bias_noise}
x_j = \mu_j + s_j,
\qquad
\mu\in\mathbb{R}^H \text{ deterministic},
\qquad
\mathbb{E}[s_j]=0.
\end{equation}

\subsection{Assumptions}

\begin{assumption}[Sub-Gaussian centered variation]
\label{assump:subg_C}
For each $j$, the centered noise $s_j$ is sub-Gaussian with parameter $\sigma^2$:
\[
\mathbb{P}(|s_j|\ge t)\le 2\exp\!\left(-\frac{t^2}{2\sigma^2}\right).
\]
The $s_j$ are independent across coordinates within a block.
\end{assumption}

The independence assumption is used only to obtain clean max bounds; weaker
dependence conditions yield similar results with different constants.

\subsection{Block scale with bias}

Define
\[
m_B \triangleq \max_{j\in B} |\mu_j|.
\]

\begin{lemma}[Upper bound]
\label{lem:block_upper}
For any $t\ge 0$, with probability at least $1-2b\exp(-t^2/(2\sigma^2))$,
\[
\alpha(x;B) \le m_B + t.
\]
\end{lemma}

\begin{proof}
For any $j\in B$, $|x_j| \le |\mu_j| + |s_j|$.
Hence
\[
\max_{j\in B}|x_j|
\le
\max_{j\in B}|\mu_j| + \max_{j\in B}|s_j|
=
m_B + \max_{j\in B}|s_j|.
\]
By a union bound,
\[
\mathbb{P}\left(\max_{j\in B}|s_j|\ge t\right)
\le
\sum_{j\in B}\mathbb{P}(|s_j|\ge t)
\le
2b\exp\!\left(-\frac{t^2}{2\sigma^2}\right).
\]
\end{proof}

\begin{lemma}[Lower bound]
\label{lem:block_lower}
For any $t\ge 0$, with probability at least $1-2\exp(-t^2/(2\sigma^2))$,
\[
\alpha(x;B) \ge m_B - t.
\]
\end{lemma}

\begin{proof}
Let $j^\star\in B$ achieve $|\mu_{j^\star}|=m_B$.
Then
\[
\alpha(x;B)=\max_{j\in B}|x_j|\ge |x_{j^\star}|
\ge |\mu_{j^\star}| - |s_{j^\star}|
= m_B - |s_{j^\star}|.
\]
Thus
\[
\mathbb{P}\big(\alpha(x;B)\le m_B - t\big)
\le
\mathbb{P}(|s_{j^\star}|\ge t)
\le
2\exp\!\left(-\frac{t^2}{2\sigma^2}\right).
\]
\end{proof}

\paragraph{Consequence.}
Combining Lemmas~\ref{lem:block_upper} and \ref{lem:block_lower} shows that
block scale concentrates around $m_B$ whenever $m_B \gg \sigma\sqrt{\log b}$.

\subsection{Block scale without bias (mean removal)}

After mean removal, the centered vector is $x'_j = s_j$ and the block scale is
\[
\alpha(x';B) = \max_{j\in B}|s_j|.
\]

\begin{lemma}[Centered max bound]
\label{lem:centered_max}
For any $\delta\in(0,1)$, with probability at least $1-\delta$,
\[
\alpha(x';B) \le \sigma \sqrt{2\log(2b/\delta)}.
\]
\end{lemma}

\begin{proof}
Set $t=\sigma\sqrt{2\log(2b/\delta)}$ in Lemma~\ref{lem:block_upper} with $m_B=0$:
\[
\mathbb{P}\!\left(\max_{j\in B}|s_j|\ge t\right)
\le 2b\exp\!\left(-\frac{t^2}{2\sigma^2}\right)
= 2b\exp\!\left(-\log(2b/\delta)\right)=\delta.
\]
\end{proof}

\subsection{Dynamic-range inflation factor}

The ratio of block scales with and without bias is
\[
\frac{\alpha(x;B)}{\alpha(x';B)}.
\]

\begin{theorem}[Bias-driven dynamic-range inflation]
\label{thm:range_inflation}
Fix $\delta\in(0,1)$.
Under Assumption~\ref{assump:subg_C}, with probability at least $1-2\delta$,
\[
\frac{\alpha(x;B)}{\alpha(x';B)}
\ge
\frac{m_B - \sigma\sqrt{2\log(2/\delta)}}{\sigma\sqrt{2\log(2b/\delta)}},
\]
and
\[
\frac{\alpha(x;B)}{\alpha(x';B)}
\le
\frac{m_B + \sigma\sqrt{2\log(2b/\delta)}}{\sigma\sqrt{2\log(2b/\delta)}}.
\]
In particular, if $m_B \gg \sigma\sqrt{\log b}$, then
\[
\alpha(x;B) \approx m_B
\qquad\text{while}\qquad
\alpha(x';B) = O(\sigma\sqrt{\log b}),
\]
so the block scale is inflated by a factor on the order of $m_B / (\sigma\sqrt{\log b})$.
\end{theorem}

\begin{proof}
The upper bound follows by combining Lemma~\ref{lem:block_upper} (with probability
$1-\delta$) and Lemma~\ref{lem:centered_max} (with probability $1-\delta$).
The lower bound follows by Lemma~\ref{lem:block_lower} and Lemma~\ref{lem:centered_max}.
A union bound gives total probability at least $1-2\delta$.
\end{proof}

\subsection{Implication for quantization resolution and tail SNR}

Block quantization step size satisfies
\[
\Delta(x;B) = \alpha(x;B)/Q.
\]
Thus Theorem~\ref{thm:range_inflation} implies that bias increases the quantization
step size multiplicatively in blocks containing large-mean coordinates.
This reduces effective resolution for coordinates whose magnitude is governed by
centered variation $s_j$, i.e., the long-tail semantic component.

A crude but informative proxy for tail signal-to-noise ratio (SNR) is
\[
\mathrm{SNR}_{\text{tail}} \propto \frac{\mathbb{E}[s_j^2]}{\Delta(x;B)^2},
\]
which decreases quadratically with block scale.
Therefore, bias-driven inflation of $\alpha(x;B)$ directly suppresses tail SNR.

\paragraph{Remark (percentile-based scaling).}
In practice, scales are often set by a high percentile (e.g.\ $p99.9$) rather than
the maximum.
The same phenomenon holds: when the mean dominates the upper quantiles of $|x_j|$,
percentile-based scales concentrate around $m_B$ rather than the noise level.
Replacing the max bounds above with sub-Gaussian quantile bounds yields analogous
conclusions.

\paragraph{Summary.}
Bias inflates blockwise quantization scales whenever blocks contain coordinates
with large $|\mu_j|$, compressing the resolution allocated to tail variability.
Mean removal eliminates this inflation and restores block scales to the
noise-controlled regime $O(\sigma\sqrt{\log b})$.
This formalizes the empirical dynamic-range reductions observed in
Section~2.4C.
\newpage
\section{Algorithm: Subspace-Aware Whitening via Cholesky QR}

\textbf{Input:}
\begin{itemize}
    \item $G \in \mathbb{R}^{d \times B}$: Input batch (gradients or activations), with feature dimension $d$ and batch size $B$.
    \item $k$: Rank of the spike subspace (typically $k \approx d/100$).
    \item $\lambda$: Whitening strength scalar (default $1.0$).
\end{itemize}

\textbf{Output:}
\begin{itemize}
    \item $G_{clean} \in \mathbb{R}^{d \times B}$: The whitened matrix with mean bias and anisotropic spikes removed.
\end{itemize}

\subsection*{Phase 1: Mean Removal (Centering)}
Remove the Rank-1 common mean shared across all tokens to eliminate the ``common energy'' bias.

\begin{align}
    \mu &= \frac{1}{B} G \mathbf{1}_B & \text{(Compute mean vector)} \\
    G_c &= G - \mu \mathbf{1}_B^\top & \text{(Broadcast subtract mean)}
\end{align}

\subsection*{Phase 2: Randomized Spike Detection (Sketching)}
Use a random probe matrix to identify the dominant $k$-dimensional subspace (the frequency spike) using a randomized range finder.

\begin{align}
    \Omega &\sim \mathcal{N}(0, 1)^{B \times k} & \text{(Generate Gaussian random probes)} \\
    Y &= G_c \Omega & \text{(Sketch the range via Matrix Multiply)}
\end{align}
\textit{Note: $Y \in \mathbb{R}^{d \times k}$ captures the dominant directions.}

\subsection*{Phase 3: Cholesky QR Orthogonalization}
Compute an orthonormal basis $Q$ for the range of $Y$. We use Cholesky QR to maintain hardware efficiency by relying solely on matrix multiplications.

\begin{align}
    M &= Y^\top Y & \text{(Compute Gram Matrix, result is $k \times k$)} \\
    R &= \text{Cholesky}(M) & \text{(Decompose such that $M = R^\top R$)} \\
    Q &= Y R^{-1} & \text{(Solve for $Q$ via triangular solve)}
\end{align}
\textit{Note: $Q \in \mathbb{R}^{d \times k}$ is now an orthonormal basis for the spike subspace.}

\subsection*{Phase 4: Whitening (Projection \& Subtraction)}
Project the centered data onto the discovered spike basis $Q$ and subtract it to whiten the spectrum.

\begin{align}
    C &= Q^\top G_c & \text{(Compute projection coefficients)} \\
    G_{clean} &= G_c - \lambda (Q C) & \text{(Subtract reconstruction of the spike)}
\end{align}
\newpage
\section{Algorithm: Subspace-Aware Whitening via Standard QR}

\textbf{Input:}
\begin{itemize}
    \item $G \in \mathbb{R}^{d \times B}$: Input batch (gradients or activations), with feature dimension $d$ and batch size $B$.
    \item $k$: Rank of the spike subspace (typically $k \approx d/100$).
    \item $\lambda$: Whitening strength scalar (default $1.0$).
\end{itemize}

\textbf{Output:}
\begin{itemize}
    \item $G_{clean} \in \mathbb{R}^{d \times B}$: The whitened matrix with mean bias and anisotropic spikes removed.
\end{itemize}

\subsection*{Phase 1: Mean Removal (Centering)}
Remove the Rank-1 common mean shared across all tokens.

\begin{align}
    \mu &= \frac{1}{B} G \mathbf{1}_B & \text{(Compute mean vector)} \\
    G_c &= G - \mu \mathbf{1}_B^\top & \text{(Broadcast subtract mean)}
\end{align}

\subsection*{Phase 2: Randomized Spike Detection (Sketching)}
Use a random probe matrix to identify the dominant $k$-dimensional subspace (the frequency spike).

\begin{align}
    \Omega &\sim \mathcal{N}(0, 1)^{B \times k} & \text{(Generate Gaussian random probes)} \\
    Y &= G_c \Omega & \text{(Sketch the range via Matrix Multiply)}
\end{align}
\textit{Note: $Y \in \mathbb{R}^{d \times k}$ captures the dominant directions.}

\subsection*{Phase 3: Orthogonalization via Thin QR}
Compute an orthonormal basis $Q$ for the range of $Y$ using standard thin QR decomposition. This provides better numerical stability than Cholesky QR for ill-conditioned matrices.

\begin{align}
    [Q, R] &= \text{QR}(Y, \text{mode='economic'})
\end{align}
\textit{Note: $Q \in \mathbb{R}^{d \times k}$ is the orthonormal basis for the spike. $R$ is ignored.}

\subsection*{Phase 4: Whitening (Projection \& Subtraction)}
Project the centered data onto the discovered spike basis $Q$ and subtract it to whiten the spectrum.

\begin{align}
    C &= Q^\top G_c & \text{(Compute projection coefficients)} \\
    G_{clean} &= G_c - \lambda (Q C) & \text{(Subtract reconstruction of the spike)}
\end{align}

%\newpage
%\input{./sections/learning}


%
\section{Empirical Discovery: Mean-Dominated Geometry}
\label{sec:discovery}

\subsection{Activation Mean Is Statistically Significant}
\label{subsec:mean_significance}

We measure the noise-normalized mean magnitude
\[
R_x = \frac{\|\mu_x\|_2}{\sqrt{\mathbb{E}\|x\|_2^2 / N}}
\]
and show that activation means satisfy $R_x \gg 1$, while gradient-output means
satisfy $R_\delta \approx 1$.
This establishes that bias is intrinsic to activations and not to gradient noise.


\subsection{Outliers Are Dominated by the Mean}
\label{subsec:outlier_dominance}

We decompose activation tensors into mean, spike, and tail components using SVD,
and analyze the top $0.1\%$ largest-magnitude elements.
We show that the majority of squared magnitude within outliers is contributed
by the mean component.


\subsection{Spectral Role of Bias}
\label{subsec:spectral_role}

We analyze how the mean contributes to the spectral spike and quantify
its fraction of spike energy.
This reveals that the bias forms a dense, coherent direction embedded
within the dominant spectral subspace.


\subsection{Bias Propagation into Gradients}
\label{subsec:bias_propagation}

Using the gradient identity
\[
\mathbb{E}[\delta x^\top]
=
\mathbb{E}[\delta]\mu_x^\top
+
\mathbb{E}[\delta \tilde{x}^\top],
\]
we show that activation bias induces a coherent rank-one component
in gradient matrices even when gradient outputs themselves have negligible mean.


\section{The Blessing: Efficient Mean Removal}
\label{sec:mean_removal}

\subsection{Mean Removal as Rank-One Suppression}

We define
\[
X_c = X - \mathbf{1}\mu^\top
\]
and show that mean removal significantly reduces extreme outliers and
partially attenuates spectral spike energy.


\subsection{Element-Wise and Block-Wise Outlier Reduction}

We quantify p99.9 entry magnitude and blockwise maxima before and after
mean removal, demonstrating substantial reduction in extreme scales.


\section{Downstream Validation}
\label{sec:applications}

\subsection{Low-Bit Quantization}
\label{subsec:quantization}

We analyze how mean removal improves signal-to-noise ratio under
blockwise FP4 quantization and compare against SVD-based spike whitening.


\subsection{Optimizer Conditioning}
\label{subsec:optimizer}

We examine the effect of bias on Adam second-moment estimates
\[
v = \mathbb{E}[g \odot g]
\]
and show that mean removal reduces pathological inflation of adaptive
learning-rate scaling.


\section{Theoretical Formalization}
\label{sec:theory}

We formalize a minimal model of activation bias and semantic variability,
derive bounds on mean-dominated outlier contributions, and prove that
mean removal preserves semantic covariance up to controlled error.


\section{Discussion and Implications}
\label{sec:discussion}

We discuss:
\begin{itemize}
    \item Why bias emerges in language representations.
    \item When mean removal suffices versus when spectral whitening is needed.
    \item Broader implications for scaling, stability, and hardware efficiency.
\end{itemize}


\section{Conclusion}
\label{sec:conclusion}

We argue that mean-dominated bias is a fundamental structural component
of LLM representations and a key hidden driver of instability.
By recognizing bias as both curse and blessing, we obtain a simple,
efficient, and theoretically grounded mechanism for stabilizing
large-scale training.

\newpage


\bibliographystyle{../doc/iclr-template/master/iclr2026/iclr2026_conference}
\bibliography{references}

\newpage

\appendix
%\input{proofs}



\end{document}